% PDF 59--112；原书第 29--82 页。
\chapter{语言模型预训练}
\label{ch:pretraining}
\chaptercredit{2}
\chapterabstract{预训练是大语言模型能力的基础。借助大规模数据集，模型获得重要的语言理解和生成能力。数据集的规模和质量关系到能否充分发挥模型的潜力。有效的预训练也离不开合适的模型结构、加速方法和优化策略。本章先介绍编码器—解码器架构，说明它适用于哪些自然语言处理任务，以及有哪些局限。接着介绍注意力机制，帮助读者理解大多数大语言模型的核心架构 Transformer。然后讨论数据收集与处理，以及模型架构、预训练目标和优化策略等关系到预训练的设计要点。我们还考察 BERT、T5、GPT（第 1—3 代）和 Mixtral~8x7B 等衍生出众多变体的主要模型，对比它们的架构与训练方式。章末的教程将深入分析大语言模型架构，说明掩码模型和因果模型的区别，考察预训练模型如何产生输出，并简要介绍训练过程。}

% PDF 59--61。
\section{编码器—解码器架构}
图~\ref{fig:ch02-encoder-decoder} 所示的编码器—解码器架构，是自然语言处理的重要进展，尤其适用于机器翻译、生成式摘要（abstractive summarization）和问答等序列到序列任务（\refcite{ch02-sutskever}{Sutskever et al., 2014}）。这个框架由编码器（encoder）和解码器（decoder）两个主要部分组成。

\subsection{编码器}
首先将输入文本切分为词元，单位可以是词或子词（subword）。然后把这些单位嵌入为特征向量 $\symbf{x}_1,\ldots,\symbf{x}_T$。单向编码器在每个时刻 $t$ 使用 $\symbf{h}_{t-1}$ 和 $\symbf{x}_t$ 更新隐藏状态（hidden state）$\symbf{h}_t$。
\begin{equation}
\symbf{h}_t=f(\symbf{h}_{t-1},\symbf{x}_t).
\tag{2.1}
\end{equation}
编码器的最终状态 $\symbf{h}_t$ 称为上下文变量（context variable）或上下文向量（context vector），编码整个输入序列的信息，写为
\begin{equation}
\symbf{c}=m(\symbf{h}_1,\ldots,\symbf{h}_T),
\tag{2.2}
\end{equation}
其中 $m$ 是映射函数。最简单的情形是用最后一个隐藏状态作为上下文变量。
\begin{equation}
\symbf{c}=m(\symbf{h}_1,\ldots,\symbf{h}_T)=\symbf{h}_T.
\tag{2.3}
\end{equation}
进一步扩展架构可以得到双向编码器。此时，隐藏状态除了依赖上一个隐藏状态 $\symbf{h}_{t-1}$ 和输入 $\symbf{x}_t$，还依赖下一个状态 $\symbf{h}_{t+1}$。

\begin{figure}[!htbp]
\centering
% 重绘自原书第 30 页的图 2.1。
\begin{tikzpicture}[
  x=.96cm,font=\footnotesize\sffamily,
  flow/.style={-{Stealth[length=1.35mm]},draw=black!70,line width=.45pt},
  wire/.style={draw=black!70,line width=.45pt},
  state/.style={circle,minimum size=6.8mm,inner sep=0pt,draw=blue!32},
  vector/.style={draw=black!58,rounded corners=1pt,minimum width=4mm,minimum height=9mm,inner sep=0pt},
  coordinate label/.style={font=\footnotesize\rmfamily}]
\draw[rounded corners=2mm,dashed,black!65] (.84,1.5) rectangle (6.16,3.12);
\draw[rounded corners=2mm,black!60] (6.34,1.5) rectangle (11.55,3.12);
\node at (3.45,2.94) {编码器};
\node at (8.96,2.94) {解码器};

\foreach \i/\x in {1/1.29,2/2.38,3/3.47,4/4.56,5/5.65} {
  \node[state,fill=blue!15!gray!11] (e\i) at (\x,2.27) {};
  \node[vector] (x\i) at (\x,.39) {};
  \fill[green!22!gray!18,draw=green!40!gray] (\x,.55) circle[radius=1.13mm];
  \fill[green!22!gray!18,draw=green!40!gray] (\x,.23) circle[radius=1.13mm];
  \draw[flow] (x\i.north) -- (e\i.south);
}
\foreach \i/\j in {1/2,2/3,3/4,4/5} {
  \draw[flow] (e\i.east) -- (e\j.west);
}

\foreach \i/\x in {1/6.81,2/7.86,3/8.91,4/9.96,5/11.01} {
  \node[state,fill=yellow!22] (d\i) at (\x,2.27) {};
  \node[vector] (in\i) at (\x+.11,.43) {};
  \node[vector] (out\i) at (\x,4.03) {};
  \foreach \offset in {-.16,.16} {
    \fill[orange!20,draw=orange!55!gray] (\x+.11,.43+\offset) circle[radius=1.13mm];
    \fill[orange!20,draw=orange!55!gray] (\x,4.03+\offset) circle[radius=1.13mm];
  }
  \draw[flow] (\x,1.12) -- (d\i.south);
  \draw[flow] (in\i.north) -- (\x+.11,1.98);
  \draw[flow] (d\i.north) -- (out\i.south);
}

\draw[wire,semithick] (e5.east) -- (6.42,2.27) |- (11.01,1.12);
\node[align=right,anchor=east] at (1.1,3.62) {隐藏状态};
\node[align=right,anchor=east] at (1.02,.48) {特征\\向量};
\node[align=right,anchor=east] at (1.02,-.43) {输入};
\node at (5.85,4.86) {输出};
\foreach \x/\sym in {1.84/1,2.93/2,6.16/t} {
  \node[coordinate label] at (\x,3.62) {$h_{\sym}$};
  \draw[flow] (\x,3.44) -- (\x,2.33);
}
\node[coordinate label] at (1.29,-.43) {$x_1$};
\node[coordinate label] at (5.65,-.43) {$x_T$};
\node[coordinate label] at (6.92,-.40) {$y_1$};
\node[coordinate label] at (11.12,-.40) {$y_U$};
\node[coordinate label] at (6.81,4.86) {$y_1$};
\node[coordinate label] at (11.01,4.86) {$y_U$};
\end{tikzpicture}

\caption{用于序列到序列处理的单向编码器—解码器。}
\label{fig:ch02-encoder-decoder}
\end{figure}

\subsection{解码器}
解码器取得编码器给出的上下文向量 $\symbf{c}$ 后，开始生成输出序列 $y=(y_1,y_2,\ldots,y_U)$，其中 $U$ 可以不同于 $T$。与编码器类似，解码器在任意时刻的隐藏状态为
\begin{equation}
\symbf{s}_{t'}=g(\symbf{s}_{t-1},\symbf{y}_{t'-1},\symbf{c}).
\tag{2.4}
\end{equation}
解码器的隐藏状态传至输出层。时刻 $t'$ 的下一个词元的条件分布为
\begin{equation}
P(\symbf{y}_{t'}\mid \symbf{y}_{t'-1},\ldots,\symbf{y}_1,\symbf{c})
=\operatorname{softmax}(\symbf{s}_{t-1},\symbf{y}_{t'-1},\symbf{c}).
\tag{2.5}
\end{equation}
\transnote{式（2.4）和（2.5）保留原书第 31 页的下标。式中解码器上一状态写为 $\symbf{s}_{t-1}$，输出时间索引则用 $t'$。}

\subsection{训练与优化}
编码器—解码器模型通过监督学习（supervised learning）进行端到端（end-to-end）训练。常用的损失函数是预测输出序列与真实输出之间的类别交叉熵（categorical cross-entropy）。
\begin{equation}
\mathcal{L}=-\sum_{t=1}^{U}\log p(y_t\mid y_{t-1},\ldots,y_1,\symbf{c}).
\tag{2.6}
\end{equation}
优化模型参数时，通常采用梯度下降（gradient descent）的变体，如 Adam 或 RMSprop。

\subsection{编码器—解码器架构的问题}
如上文所述，编码器把源句子的信息压缩到单一的上下文变量 $\symbf{c}$，供解码器使用。这种压缩不可避免地造成信息损失，输入越长，问题越突出。自然语言的句法和语义结构复杂，词元之间往往存在长程依赖（long-range dependency），很难用单个上下文向量充分表示。不过，编码器各时间步的隐藏状态仍然包含可供解码器使用的有价值信息。在每个解码时间步，这些隐藏状态可以发挥不同程度的作用，由此部分缓解只使用一个上下文变量带来的限制。循环神经网络是许多编码器—解码器模型的基础架构，却存在梯度消失（vanishing gradients）和梯度爆炸（exploding gradients）等问题（\refcite{ch02-hochreiter}{Hochreiter, 1998}）。此外，循环神经网络固有的顺序依赖使并行化更加困难，也制约了计算效率。

% PDF 62--63。
\section{注意力机制}
注意力机制有助于解决基于循环神经网络的编码器—解码器架构中的问题。如图~\ref{fig:ch02-attention} 所示，可以把注意力机制看成一个存储库。收到查询时，它根据存储的键和值产生输出（\refcite{ch02-bahdanau}{Bahdanau et al., 2014}）。

\begin{figure}[!htbp]
\centering
% 重绘自原书第 32 页的图 2.2。
\begin{tikzpicture}[
  font=\footnotesize\sffamily,
  flow/.style={draw=black!60,line width=.45pt,-{Stealth[length=1.1mm]}},
  wire/.style={draw=black!60,line width=.45pt},
  value/.style={draw=black!65,rounded corners=1pt,fill=white,minimum width=3.6mm,minimum height=7.3mm,inner sep=0pt},
  op/.style={circle,draw=black!55,fill=white,minimum size=3.7mm,inner sep=0pt}]
\filldraw[fill=yellow!18,draw=black!55,rounded corners=4mm] (2.48,-.46) rectangle (7.33,4.01);
\node[font=\footnotesize\sffamily\bfseries] at (4.94,3.78) {存储库};
\node at (2.95,3.53) {值};
\node at (2.95,.14) {键};

\node[op,minimum size=4.5mm] (sum) at (8.25,2.33) {$\sum$};
\node[value,minimum width=8.8mm,minimum height=4mm] (output) at (8.25,3.27) {};
\node[above=1pt of output] {输出};
\fill[orange!55!gray!20] (8.15,3.27) circle[radius=1mm];
\fill[orange!55!gray!20] (8.35,3.27) circle[radius=1mm];
\draw[flow] (sum.north) -- (output.south);
\node[anchor=west] at (8.56,2.33) {$\symbf{o}=\sum_{i=1}^{n}b_i\symbf{v}_i$};

\node[value,minimum width=8.8mm,minimum height=4mm] (query) at (8.25,.61) {};
\node[below=1pt of query] {查询 $\symbf{q}$};
\fill[violet!68] (8.15,.61) circle[radius=1mm];
\fill[violet!68] (8.35,.61) circle[radius=1mm];
\draw[wire] (query.west) -- (6.76,.61) -- (6.76,1.16) -- (3.68,1.16);

\foreach \x/\i/\delta in {3.82/1/.44,4.99/2/.27,6.16/n/.10} {
  \node[value] (v\i) at (\x,3.06) {};
  \node[value] (k\i) at (\x,.48) {};
  \node[op] (m\i) at (\x,2.33) {$\times$};
  \fill[green!34!gray!25] (\x,3.17) circle[radius=.94mm];
  \fill[green!34!gray!25] (\x,2.96) circle[radius=.94mm];
  \fill[blue!60!gray!35] (\x,.58) circle[radius=.94mm];
  \fill[blue!60!gray!35] (\x,.38) circle[radius=.94mm];
  \fill[blue!43!gray!28] (\x,1.53) circle[radius=1.15mm];
  \fill[blue!43!gray!28] (\x,1.91) circle[radius=1.15mm];
  \node[anchor=west] at (\x+.13,1.53) {$a_{\i}$};
  \node[anchor=west] at (\x+.13,1.91) {$b_{\i}$};
  \node[above=1pt of v\i] {$\symbf{v}_{\i}$};
  \node[below=1pt of k\i] {$\symbf{k}_{\i}$};
  \draw[flow] (k\i.north) -- (\x,1.41);
  \draw[flow] (\x-.14,1.16) -- (\x-.14,1.36) -- (\x-.09,1.46);
  \draw[flow] (\x,1.64) -- (\x,1.80);
  \draw[flow] (\x,2.02) -- (m\i.south);
  \draw[flow] (v\i.south) -- (m\i.north);
  \draw[flow] (m\i.east) -- (\x+.27,2.33) -- (\x+.27,2.33+\delta)
    -- (7.76,2.33+\delta) -- (sum.west);
}
\node[anchor=east] at (2.36,1.91) {$b_i=\dfrac{\exp(a_i)}{\sum_j\exp(a_j)}$};
\node[anchor=east] at (2.36,1.30) {$a_i=\alpha(\symbf{q},\symbf{k}_i)$};
\end{tikzpicture}

\caption{注意力机制中查询、键、值和输出向量之间的关系。}
\label{fig:ch02-attention}
\end{figure}

设存储单元由 $n$ 个键值对（key-value pair）$(\symbf{k}_1,\symbf{v}_1),\ldots,(\symbf{k}_n,\symbf{v}_n)$ 组成，其中 $\symbf{k}_i\in\mathbb{R}^{d_k}$，$\symbf{v}_i\in\mathbb{R}^{d_v}$。注意力层接收查询（query）$\symbf{q}\in\mathbb{R}^{d_q}$，返回与值 $\symbf{v}$ 形状相同的输出 $\symbf{o}\in\mathbb{R}^{d_v}$。

注意力层使用评分函数（score function）$\alpha$ 度量查询和键的相似程度。对键 $\symbf{k}_1,\ldots,\symbf{k}_n$，相应分数 $a_1,\ldots,a_n$ 由下式给出。
\begin{equation}
a_i=\alpha(\symbf{q},\symbf{k}_i).
\tag{2.7}
\end{equation}
对分数应用 softmax 函数得到注意力权重。
\begin{equation}
\symbf{b}=\operatorname{softmax}(\symbf{a}).
\tag{2.8}
\end{equation}
$\symbf{b}$ 的各个分量为
\begin{equation}
b_i=\frac{\exp(a_i)}{\sum_j\exp(a_j)}.
\tag{2.9}
\end{equation}
输出是用注意力权重对值作加权求和的结果。
\begin{equation}
\symbf{o}=\sum_{i=1}^{n} b_i\symbf{v}_i.
\tag{2.10}
\end{equation}
评分函数 $\alpha(\symbf{q},\symbf{k})$ 有多种形式，因而产生了多种注意力机制。最简单的是基于点积的评分函数，无需可调参数。一种变体是缩放点积（scaled dot product），即用 $\sqrt{d_k}$ 归一化点积，以减轻维数增大的影响（\refcite{ch02-luong}{Luong et al., 2015}；\refcite{ch02-vaswani}{Vaswani et al., 2017}）。
\begin{equation}
\alpha(\symbf{q},\symbf{k})=\frac{\symbf{q}\cdot\symbf{k}}{\sqrt{d_k}}.
\tag{2.11}
\end{equation}

\subsection{自注意力}
在自注意力（self-attention）中，每个输入向量 $\symbf{x}_i$ 都投影为三个不同的向量，分别是查询 $\symbf{q}_i$、键 $\symbf{k}_i$ 和值 $\symbf{v}_i$。用可学习的权重矩阵 $\symbf{W}_Q$、$\symbf{W}_K$ 和 $\symbf{W}_V$ 进行这些投影，分别得到 $\symbf{q}_i=\symbf{x}_i\symbf{W}_q$、$\symbf{k}_i=\symbf{x}_i\symbf{W}_k$ 和 $\symbf{v}_i=\symbf{x}_i\symbf{W}_v$。这些权重矩阵随机初始化，并在训练中优化。把查询、键和值各自组成矩阵后，可用一次矩阵计算写出这个过程。
\begin{equation}
\operatorname{attention}(\symbf{Q},\symbf{K},\symbf{V})
=\operatorname{softmax}\left(\frac{\symbf{Q}\symbf{K}^{\mathrm{T}}}{\sqrt{d_k}}\right)\symbf{V}.
\tag{2.12}
\end{equation}

% PDF 63--73；原书第 33--43 页。
\section{Transformer}
\label{sec:ch02-transformer}
\refcite{ch02-vaswani}{Vaswani 等人（2017）}提出的 Transformer 模型，是序列到序列任务的基础架构之一。图~\ref{fig:ch02-transformer} 所示架构采用编码器—解码器结构。编码器与解码器各由若干结构相同的层组成，下面介绍其中的主要部件。

\subsection{编码器}
编码器负责处理输入序列，把信息压缩为供解码器使用的上下文或记忆（memory）。每个编码器层由三个主要部件组成。
\begin{itemize}
\item \textbf{多头注意力（multi-head attention）。}每个注意力头可以关注输入的不同部分，从而捕捉数据的多个方面。
\item \textbf{前馈神经网络。}这是一个简单而有效的神经网络。它对注意力向量施加非线性变换，把结果传给下一个编码器层，也可供解码器层使用。
\item \textbf{相加与归一化（Add \& Norm）。}该层结合残差连接（residual connection）和层归一化（layer normalization），使激活值更平稳，改善训练过程，并缓解编码器和解码器中的梯度消失问题。
\end{itemize}

\begin{figure}[!htbp]
\centering
% 重绘自原书第 34 页的图 2.3。
\begin{tikzpicture}[
  font=\footnotesize\sffamily,
  layer/.style={draw=black!40,rounded corners=1.5pt,minimum width=26mm,minimum height=6.5mm,align=center,inner sep=2pt},
  norm/.style={layer,fill=yellow!37},
  ff/.style={layer,fill=cyan!19},
  attn/.style={layer,fill=orange!25,minimum height=8mm},
  embed/.style={layer,fill=red!20,minimum height=8mm},
  plus/.style={circle,minimum size=4mm,draw=black!60,inner sep=0pt},
  flow/.style={draw=black!63,line width=.56pt,-{Stealth[length=1.6mm]}},
  wire/.style={draw=black!63,line width=.56pt}]
\filldraw[fill=black!5,draw=black!20,rounded corners=3mm] (-1.56,1.45) rectangle (1.57,5.41);
\filldraw[fill=black!5,draw=black!20,rounded corners=3mm] (2.53,1.45) rectangle (5.66,7.87);
\node at (-1.91,3.55) {$N\times$};
\node at (6.05,4.67) {$N\times$};

\node (einput) at (0,-.50) {输入};
\node[embed] (eembed) at (0,.24) {输入嵌入};
\node[plus] (epos) at (0,1.16) {$+$};
\node[plus,minimum size=6mm] (ewave) at (-.82,1.16) {$\sim$};
\node[anchor=east,align=center] at (-1.16,1.16) {位置\\编码};
\node[attn] (eattn) at (0,2.20) {多头注意力};
\node[norm] (enorm1) at (0,3.12) {相加与归一化};
\node[ff] (eff) at (0,4.10) {前馈网络};
\node[norm] (enorm2) at (0,5.02) {相加与归一化};

\node[align=center] (dinput) at (4.10,-.50) {输出（右移）};
\node[embed] (dembed) at (4.10,.24) {输出嵌入};
\node[plus] (dpos) at (4.10,1.16) {$+$};
\node[plus,minimum size=6mm] (dwave) at (4.92,1.16) {$\sim$};
\node[anchor=west,align=center] at (5.26,1.16) {位置\\编码};
\node[attn] (dattn1) at (4.10,2.20) {掩码多头注意力};
\node[norm] (dnorm1) at (4.10,3.12) {相加与归一化};
\node[attn] (dattn2) at (4.10,4.28) {多头注意力};
\node[norm] (dnorm2) at (4.10,5.22) {相加与归一化};
\node[ff] (dff) at (4.10,6.24) {前馈网络};
\node[norm] (dnorm3) at (4.10,7.20) {相加与归一化};
\node[layer,fill=blue!17!gray!12] (linear) at (4.10,8.30) {线性变换};
\node[layer,fill=green!21] (softmax) at (4.10,9.23) {Softmax};
\node (output) at (4.10,10.02) {输出概率};

\foreach \origin/\target in {einput/eembed,eembed/epos,epos/eattn,eattn/enorm1,enorm1/eff,eff/enorm2,
 dinput/dembed,dembed/dpos,dpos/dattn1,dattn1/dnorm1,dattn2/dnorm2,
 dnorm2/dff,dff/dnorm3,dnorm3/linear,linear/softmax,softmax/output} {
  \draw[flow] (\origin.north) -- (\target.south);
}
\draw[wire] (ewave.east) -- (epos.west);
\draw[wire] (dwave.west) -- (dpos.east);
\foreach \x in {0,4.10} {
  \foreach \off in {-.62,.62} {
    \draw[flow] (\x,1.63) -- (\x+\off,1.63) -- (\x+\off,1.80);
  }
}
\draw[flow] (0,1.56) -- (-1.45,1.56) |- (enorm1.west);
\draw[flow] (0,3.58) -- (-1.45,3.58) |- (enorm2.west);
\draw[flow] (4.10,1.56) -- (5.54,1.56) |- (dnorm1.east);
\draw[flow] (4.10,3.55) -- (5.54,3.55) |- (dnorm2.east);
\draw[flow] (4.10,5.71) -- (5.54,5.71) |- (dnorm3.east);
\draw[wire] (enorm2.north) -- (0,5.54) -- (2.17,5.54) -- (2.17,3.64) -- (4.10,3.64);
\draw[flow] (3.48,3.64) -- (3.48,3.88);
\draw[flow] (4.10,3.64) -- (4.10,3.88);
\draw[flow] (dnorm1.north) -- (4.10,3.49) -- (4.72,3.49) -- (4.72,3.88);
\end{tikzpicture}

\caption{Transformer 架构包含编码器和解码器，两者都使用多头注意力。}
\label{fig:ch02-transformer}
\end{figure}

\subsection{解码器}
解码器接收编码器给出的上下文，并生成输出序列。它也由多层组成，与编码器有许多共同之处，但有几处改动。
\begin{itemize}
\item \textbf{掩码多头注意力（masked multi-head attention）。}与多头注意力类似，不过使用掩码机制，以确保对当前词的预测不依赖于序列中还未出现的词。
\item \textbf{编码器—解码器注意力（encoder-decoder attention）。}该层利用编码器提供的上下文，使解码器能够关注输入序列中的相关部分。
\item \textbf{前馈神经网络。}该层的架构与编码器中的相同，对注意力向量进一步变换，为生成输出序列作准备。
\end{itemize}

下面介绍 Transformer 架构的各个部件及其组成。

\subsection{词元切分与表示}
在 Transformer 模型中，词元切分通常把句子转换为机器可读的形式。根据具体应用需要的粒度，可以按词或按子词切分。按词切分时，句子中的每个词作为一个独立词元，然后映射为对应的向量表示，如词嵌入，用作 Transformer 的输入。这种办法处理词表外词（out-of-vocabulary word）时可能受到限制。字节对编码（Byte-Pair Encoding，BPE）或 WordPiece 等子词方法，通常能弥补按词切分的不足。这些方法把词分成更小的片段，即子词，使词表外词能够得到表示，也能捕捉词形上的细微差别。然后把各子词映射为嵌入并送入 Transformer。

例如，单词“unhappiness”可以切成“un”和“happiness”等子词。再分别把这些子词映射为各自的嵌入。这种方法可以提高模型的泛化能力，让模型处理更广泛的词汇，包括训练中未出现的词。

也可以采用结合按词和按子词切分的混合方法，兼顾子词表示的覆盖范围和词级词元的可解释性。

\subsection{位置编码}
Transformer 模型并行处理输入序列中的所有词元，模型本身没有考虑词元位置或顺序的机制。位置编码（positional encoding）为模型提供有关词元在序列中相对位置的信息。输入嵌入送入 Transformer 前，通常会先加上位置编码。

若句子长度为 $l$，嵌入的维数或深度为 $d$，则位置编码 $\symbf{P}$ 是一个形状相同的二维矩阵，即 $\symbf{P}\in\mathbb{R}^{l\times d}$。对于每个位置，沿长度方向的索引为 $i$，沿嵌入维度方向的索引为 $j$，可写为
\begin{equation}
\symbf{P}_{i,2j}=\sin(i/1000^{2j/d}),
\tag{2.13}
\end{equation}
\begin{equation}
\symbf{P}_{i,2j+1}=\cos(i/1000^{2j/d}),
\tag{2.14}
\end{equation}
其中 $i=0,\ldots,l-1$，$j=0,\ldots,\lfloor(d-1)/2\rfloor$。从函数定义可知，沿向量维度方向频率逐渐降低，波长从 $2\pi$ 到 $10000\cdot 2\pi$ 成等比数列。图~\ref{fig:ch02-position} 可视化了维数 $d=512$、最大位置长度 $l=100$ 时的位置编码。
\transnote{式（2.13）和（2.14）中的底数 $1000$ 及原书下一页文字中的 $10000$ 均按印刷版保留。}

\sourcefigure{fig02_04.pdf}{100 个位置、512 个维度的位置编码。横轴为嵌入维度索引，纵轴为位置索引。}{fig:ch02-position}

\subsection{多头注意力}
多头注意力使用 $h$ 个并行的自注意力头，取代单个自注意力头，以增强模型的表示能力。原始 Transformer 使用 $h=8$ 个头，使模型能够捕捉输入数据的不同方面和依赖关系，例如机器翻译任务中的语法与时态。

在多头注意力中，每个头都有自己的一组可学习的查询、键和值的权重矩阵，从而得到各自的查询、键和值矩阵及输出矩阵。把这些输出矩阵拼接起来，再用一个权重矩阵对拼接结果作线性变换。各个头从输入到输出的并行变换见图~\ref{fig:ch02-multihead}。
\begin{equation}
\mathrm{head}_i
=\operatorname{attention}(\symbf{W}_{Q_i}\symbf{Q},\symbf{W}_{K_i}\symbf{K},\symbf{W}_{V_i}\symbf{V}).
\tag{2.15}
\end{equation}
\begin{equation}
\operatorname{multihead}(\symbf{Q},\symbf{K},\symbf{V})
=\symbf{W}_O\operatorname{concat}(\mathrm{head}_1,\ldots,\mathrm{head}_h).
\tag{2.16}
\end{equation}

\begin{figure}[!htbp]
\centering
% 重绘自原书第 37 页的图 2.5。
\begin{tikzpicture}[
  font=\small\sffamily,
  unit/.style={draw=black!38,rounded corners=1.4pt,minimum height=5.9mm,minimum width=15.5mm,align=center},
  wire/.style={-{Stealth[length=1.5mm]},draw=black!62,line width=.52pt}]
\foreach \offset in {.27,.14} {
  \foreach \x in {-1.75,0,1.75} {
    \node[unit,draw=green!35!gray!20,fill=green!12] at (\x+\offset,.95+\offset) {};
    \draw[wire,opacity=.23] (\x+\offset,1.23+\offset) -- (\x+\offset,1.79+\offset);
  }
  \node[unit,minimum width=53mm,minimum height=10mm,fill=blue!12,draw=blue!18] at (\offset,2.30+\offset) {};
  \draw[wire,opacity=.24] (\offset,2.80+\offset) -- (\offset,3.39);
}
\foreach \x/\symbol in {-1.75/V,0/K,1.75/Q} {
  \node (input\symbol) at (\x,.15) {$\symbf{\symbol}$};
  \node[unit,fill=green!17] (proj\symbol) at (\x,.95) {线性变换};
  \draw[wire] (input\symbol.north) -- (proj\symbol.south);
  \draw[wire] (proj\symbol.north) -- (\x,1.79);
}
\node[unit,minimum width=53mm,minimum height=10mm,fill=blue!22!violet!12] (attention) at (0,2.30) {缩放点积注意力};
\draw[black!45] (2.99,2.27) -- (3.10,2.27) -- (3.10,3.02) -- (2.99,3.02);
\node[anchor=west] at (3.12,2.65) {$h$};
\node[unit,fill=yellow!32] (concat) at (0,3.67) {拼接};
\node[unit,fill=green!17] (linear) at (0,4.56) {线性变换};
\draw[wire] (attention.north) -- (concat.south);
\draw[wire] (concat.north) -- (linear.south);
\draw[wire] (linear.north) -- (0,5.23);
\end{tikzpicture}

\caption{多头注意力机制学习多组查询、键和值的变换矩阵，再把各头结果合并。}
\label{fig:ch02-multihead}
\end{figure}

\subsection{逐位置前馈神经网络}
注意力机制之后，Transformer 架构的下一部件是前馈神经网络。它进一步变换注意力向量，使其适合作为后续编码器层或解码器层的输入。前馈神经网络通常包含两层，在第一层之后使用整流线性单元（Rectified Linear Unit，ReLU）激活函数，引入非线性。设 $\symbf{z}$ 为输入的注意力向量，前馈神经网络施加的变换 $F(\symbf{z})$ 可表示为
\begin{equation}
F(\symbf{z})=\operatorname{ReLU}(\symbf{z}\symbf{W}_1+\symbf{b}_1)\symbf{W}_2+\symbf{b}_2,
\tag{2.17}
\end{equation}
其中 $\symbf{W}_1$、$\symbf{W}_2$ 分别是第一、二层的权重矩阵，$\symbf{b}_1$、$\symbf{b}_2$ 分别是两层的偏置向量。Transformer 编码器和解码器的 $N$ 层中，每层对输入序列的各个词使用相同的线性变换，各层的权重参数 $(\symbf{W}_1,\symbf{W}_2)$ 和偏置参数 $(\symbf{b}_1,\symbf{b}_2)$ 则各自独立。

\subsection{层归一化}
Transformer 与残差网络（Residual Network，ResNet）类似，也使用残差连接，将输入 $\symbf{X}$ 加到输出 $\symbf{Z}$ 上（\refcite{ch02-he}{He et al., 2016}）。归一化过程使每层的激活值具有零均值和单位方差。对于每个隐藏单元 $h_i$，层归一化写作
\begin{equation}
h_i=\frac{g}{\sigma}(h_i-\mu),
\tag{2.18}
\end{equation}
其中 $g$ 是增益变量（通常设为 1），均值 $\mu=\frac{1}{H}\sum_{i=1}^{H}h_i$，标准差 $\sigma=\sqrt{\frac{1}{H}\sum_{i=1}^{H}(h_i-\mu)^2}$。

层归一化技术减小协变量偏移（covariate shift）。这里指层之间的梯度依赖。所需的迭代次数因而减少，收敛加快（\refcite{ch02-ba}{Ba et al., 2016}）。

\subsection{掩码多头注意力}
在 Transformer 中，解码器根据编码器的输出以及目标序列中已经见到的词元，预测序列中的下一个词元（词或字符）。解码器第一层采用一种特殊策略，只能访问在当前待预测词元之前的词元。这种机制称为掩码多头注意力。

掩码用一个特定的权重矩阵 $\symbf{M}$ 实现。其中，与序列未来词元对应的元素设为 $-\infty$，与此前词元对应的元素设为 0。

在计算查询矩阵 $\symbf{Q}$ 与键矩阵的转置 $\symbf{K}^{\mathrm{T}}$ 的乘积之后、应用 softmax 之前施加掩码。于是未来词元对应的 softmax 输出为零，从计算中被屏蔽。这样，解码器便无法预先访问序列中的未来词元，能够保持机器翻译等任务所要求的序列顺序。
\begin{equation}
\operatorname{maskedAttention}(\symbf{Q},\symbf{K},\symbf{V})
=\operatorname{softmax}\left(\frac{\symbf{Q}\symbf{K}^{\mathrm{T}}+\symbf{M}}{\sqrt{d_k}}\right)\symbf{V}.
\tag{2.19}
\end{equation}

\subsection{编码器—解码器注意力}
编码器—解码器注意力机制连接编码器与解码器，把上下文信息从源序列传给目标序列。从概念上说，这一层的工作方式与标准多头注意力相似。关键区别在于，查询 $\symbf{Q}$ 来自解码器的当前状态，而键 $\symbf{K}$ 和值 $\symbf{V}$ 取自编码器的输出。这样，生成目标序列中的每个词元时，模型可以关注源序列的相关部分，从而捕捉源序列与目标序列之间复杂的关系。

\subsection{Transformer 变体}
研究者对上述原始 Transformer 作了多种改动，产生了许多 Transformer 模型（\refcite{ch02-lin}{Lin et al., 2022}）。如图~\ref{fig:ch02-taxonomy} 所示，可以把这些改动分为三类，即架构改动、预训练方法和应用。下面各节详细介绍不同变体之间的主要差别，表~\ref{tab:ch02-variants} 列出了部分变体。

\begin{figure}[!htbp]
\centering
% 重绘自原书第 39 页的图 2.6。
\begin{tikzpicture}[
  font=\scriptsize\sffamily,
  block/.style={draw=black!28,rounded corners=1.2pt,align=center,
     minimum height=4.9mm,minimum width=18mm,inner sep=1.4pt},
  items/.style={rounded corners=1.3pt,draw=black!25,text width=26mm,
     align=left,inner sep=2.8pt,font=\scriptsize\sffamily},
  frame/.style={dashed,draw=black!45,rounded corners=2pt,line width=.54pt},
  wire/.style={-{Stealth[length=1.2mm]},draw=black!62,line width=.47pt}]
\draw[frame] (0,-.23) rectangle (3.00,6.37);
\node[anchor=west,font=\scriptsize\sffamily\bfseries] at (.13,6.12) {模块层面};
\node[items,fill=cyan!15,anchor=north] (act) at (1.50,5.88) {激活函数\\扩大容量\\去掉 FFN 模块};
\node[items,fill=yellow!31,below=1.5mm of act] (norm) {归一化的位置\\替代归一化\\免归一化};
\node[items,fill=red!22,below=1.5mm of norm] (attention) {稀疏注意力\\线性化注意力\\查询原型\\记忆压缩\\低秩自注意力\\带先验的注意力\\改进的多头注意力};
\node[items,fill=blue!19,below=1.5mm of attention] {绝对位置\\相对位置\\其他表示\\隐式表示};

\filldraw[fill=black!5,draw=black!20,rounded corners=2mm] (4.20,2.01) rectangle (6.27,6.07);
\node[block,fill=yellow!28] (n2) at (5.23,5.73) {相加与归一化};
\node[block,fill=cyan!17,minimum height=7mm] (ff) at (5.23,4.96) {逐位置 FFN};
\node[block,fill=yellow!28] (n1) at (5.23,3.97) {相加与归一化};
\node[block,fill=red!20,minimum height=7mm] (ma) at (5.23,3.20) {多头注意力};
\node[draw=black!56,circle,inner sep=1pt] (p) at (5.23,1.63) {$+$};
\node[block,fill=blue!16,minimum width=14mm] (pos) at (4.14,1.63) {位置编码};
\node[block,fill=cyan!17] (token) at (5.23,.79) {词元嵌入};
\node (input) at (5.23,.11) {输入};
\node at (6.54,4.48) {$\times L$};
\draw[wire] (input.north) -- (token.south);
\draw[wire] (token.north) -- (p.south);
\draw[wire] (pos.east) -- (p.west);
\draw[wire] (p.north) -- (ma.south);
\foreach \offset in {-.43,.43} {
  \draw[wire] (5.23,2.54) -- (5.23+\offset,2.54) -- (5.23+\offset,2.85);
}
\draw[wire] (ma.north) -- (n1.south);
\draw[wire] (n1.north) -- (ff.south);
\draw[wire] (ff.north) -- (n2.south);
\draw[wire] (n2.north) -- (5.23,6.66);
\draw[wire] (5.23,2.18) -- (4.28,2.18) |- (n1.west);
\draw[wire] (5.23,4.37) -- (4.28,4.37) |- (n2.west);

\draw[frame] (7.03,3.59) rectangle (11.56,5.99);
\node[anchor=west,font=\scriptsize\sffamily\bfseries] at (7.17,5.72) {架构层面};
\node[items,fill=green!16,text width=39mm,anchor=north] at (9.31,5.51)
  {轻量级变体\\跨块连接\\自适应计算时间\\循环与层次结构\\替代架构};
\draw[black!38,densely dotted] (6.28,5.77) -- (7.03,5.95);
\draw[black!38,densely dotted] (6.28,2.20) -- (7.03,3.62);

\draw[frame] (7.03,2.35) rectangle (11.56,3.44);
\node[anchor=west,font=\scriptsize\sffamily\bfseries] at (7.17,3.23) {预训练模型};
\foreach \x/\name in {7.85/编码器,9.30/解码器,10.77/{编码器\\—解码器}} {
  \node[block,fill=cyan!15,minimum width=13mm,minimum height=5.3mm] at (\x,2.70) {\name};
}
\draw[frame] (7.03,1.22) rectangle (11.56,2.20);
\node[anchor=west,font=\scriptsize\sffamily\bfseries] at (7.17,2.00) {应用};
\foreach \x/\name in {7.70/文本,8.76/视觉,9.82/音频,10.88/多模态} {
  \node[block,fill=orange!22,minimum width=9.5mm] at (\x,1.55) {\name};
}
\end{tikzpicture}

\caption{Transformer 已衍生出许多变体，可按架构改动、训练目标和输入类型来分类。}
\label{fig:ch02-taxonomy}
\end{figure}

\paragraph{归一化方法}
大语言模型在预训练阶段容易出现训练不稳定的问题。归一化方法有助于平稳训练。早期常用批归一化（Batch Normalization，BatchNorm），但在序列长度可变、批量较小的情形下效率不佳。因此引入了层归一化（Layer Normalization，LayerNorm 或 LN），按层归一化激活值，分别为各层重新计算均值和方差。后来提出均方根归一化（Root Mean Square Layer Normalization，RMSNorm），根据激活值的均方根重新缩放激活值，以加快层归一化的训练速度。这种方法在 Transformer 中展现出更快的训练速度和更好的表现。

\begin{practicaltip}
Gopher 和 Chinchilla 等模型采用了 RMSNorm。Microsoft 提出的 DeepNorm 能帮助深层 Transformer 在训练时保持平稳，使其可扩展到 1\,000 层。GLM-130B 等对训练稳定性和性能有要求的模型采用了这一方法。
\end{practicaltip}

\paragraph{归一化位置}
归一化主要有三种放置方式，即后置层归一化（post-LN）、前置层归一化（pre-LN）和夹层式层归一化（sandwich-LN）。原始 Transformer 采用后置归一化，放在残差块之间。不过，由于输出层附近的梯度较大，采用后置归一化的 Transformer 往往在训练中不够平稳。因此，后置归一化在大语言模型中较少单独使用，但可以配合其他策略，例如在 GLM-130B 中与前置归一化结合。前置归一化放在各子层之前，在最终预测之前还额外作一次层归一化。

\begin{practicaltip}
采用前置归一化的 Transformer 比采用后置归一化的模型更容易平稳训练，但性能有所下降。尽管有这一取舍，前置归一化通常因训练较平稳而被采用。也有例外，例如参数量超过 1\,000 亿的 GLM 等模型采用前置归一化时，训练并不平稳。夹层式层归一化是前置归一化的扩展，在残差连接前额外加入一层 LN，以缓解 Transformer 层输出值爆炸。不过，这种做法并非总能使大语言模型的训练趋于平稳，也可能导致训练崩溃。
\end{practicaltip}

\paragraph{激活函数}
恰当地选择激活函数，对于发挥前馈网络的性能很重要。已有的大语言模型通常采用高斯误差线性单元（Gaussian Error Linear Unit，GeLU）。

\begin{practicaltip}
PaLM、LaMDA 等较新的大语言模型采用门控线性单元（Gated Linear Unit，GLU）的变体，包括 SwiGLU 和 GeGLU，通常能改善表现。不过，与 GeLU 相比，这些变体所用的前馈网络参数约多 50\%。
\end{practicaltip}

\paragraph{位置嵌入}
原始 Transformer 采用绝对位置嵌入（absolute position embedding），在编码器和解码器底部加到输入嵌入上。它有正弦式和学习式两种变体，后者在预训练语言模型中较为常见。

\begin{practicaltip}
相对位置嵌入（relative position embedding）根据键与查询之间的偏移生成，最早在 Transformer-XL 中引入，后来在 T5 中作了修改。T5 根据查询与键的位置距离，在注意力分数上加一个可学习的标量，以简化嵌入。采用相对位置嵌入的 Transformer 能处理比训练时所见序列更长的序列。旋转位置嵌入（Rotary Position Embedding，RoPE）使用取决于词元位置的旋转矩阵，从而能计算带有相对位置信息的分数。由于表现良好并具有长距离衰减性质，RoPE 被用于 PaLM 和 Llama 等较新的大语言模型。xPos 基于 RoPE，在旋转角度向量的各个维度引入特定的指数衰减，以增强 Transformer 的平移不变性（translation invariance）和长度外推（length extrapolation）能力，使训练在较远距离上更平稳。ALiBi 用取决于键与查询之间距离的惩罚项调整注意力分数，改善 Transformer 的外推能力，且无需可训练参数。与正弦位置编码、RoPE、T5 位置偏置等其他位置嵌入方法相比，它展现出更好的外推性能和训练稳定性。
\end{practicaltip}

\paragraph{注意力机制}
原始 Transformer 使用全注意力（full attention），也就是对序列中的每一对词元计算注意力。它采用缩放点积注意力和多头注意力。各头对查询、键和值作不同的投影，然后拼接各头的输出，得到最终输出。全注意力的计算复杂度随序列长度呈二次增长，长序列的问题尤其突出，稀疏注意力（sparse attention）可以缓解这一问题。

\begin{practicaltip}
采用局部带状稀疏注意力（locally banded sparse attention）的高效 Transformer 变体，例如 GPT-3 的分解注意力（Factorized Attention），让每个查询只关注由位置确定的词元子集，降低计算复杂度。多查询注意力（multi-query attention）让各头共享键和值的线性变换矩阵，既节约计算，模型质量又几乎不受影响。PaLM 和 StarCoder 等模型采用了多查询注意力。FlashAttention 在不影响模型质量的前提下优化了 GPU 上注意力模块的速度和显存占用。它把输入分块并引入重计算，以便高效利用 GPU 上的快速存储器，即静态随机存取存储器（Static Random Access Memory，SRAM）。FlashAttention 集成在 PyTorch、DeepSpeed 和 Megatron-LM 等平台中，从考虑输入输出（I/O）开销的角度优化注意力模块。为了兼顾泛化和训练稳定性，建议用前置 RMSNorm（pre-RMSNorm）作层归一化，并用 SwiGLU 或 GeGLU 作激活函数。为避免性能下降，建议避免紧接在嵌入层之后使用层归一化。Realformer 和 Predictive Attention Transformer 等方法复用此前块的注意力分布来引导当前块，从而在网络中形成更直接的传播路径。透明注意力（Transparent Attention）在交叉注意力（cross-attention）模块中使用所有编码器层表示的加权和，使优化更容易。自适应计算时间（Adaptive Computation Time，ACT）根据输入的难度调整计算时间，衍生出 Universal Transformer 和 Conditional Computation Transformer 等方法。它们迭代地改进表示，或者利用门控机制优化计算资源。
\end{practicaltip}

% 按原书第 42 页重排表 2.1。
\begin{table}[!htbp]
\centering
\caption{Transformer 的网络配置。$\operatorname{Sublayer}$ 指 Transformer 层中的前馈神经网络（FFN）或自注意力模块。$d$ 表示网络隐藏状态的维数，$p_i$ 表示位置 $i$ 的嵌入。$A_{ij}$ 是给定查询与其对应键之间的注意力分数，$r_{i-j}$ 是表示查询与键之间位置差的可学习标量。$R_{\theta,t}$ 是旋转矩阵，旋转角度为 $t$ 与 $\theta$ 的乘积。}
\label{tab:ch02-variants}
\footnotesize
\setlength{\tabcolsep}{4pt}
\renewcommand{\arraystretch}{1.23}
\begin{tabularx}{\linewidth}{@{}p{.18\linewidth}p{.21\linewidth}>{\raggedright\arraybackslash}X@{}}
\toprule
配置 & 方法 & 公式\\
\midrule
归一化位置 & 后置归一化 [1] & $\operatorname{Norm}(\symbf{x}+\operatorname{Sublayer}(\symbf{x}))$\\
 & 前置归一化 [2] & $\symbf{x}+\operatorname{Sublayer}(\operatorname{Norm}(\symbf{x}))$\\
 & 夹层式归一化 [3] & $\symbf{x}+\operatorname{Norm}(\operatorname{Sublayer}(\operatorname{Norm}(\symbf{x})))$\\
\midrule
归一化方法 & LayerNorm [4] &
$\begin{aligned}[t]
 & \tfrac{\symbf{x}-\mu}{\sqrt{\sigma}}\cdot\gamma+\beta,\quad
 \mu=\tfrac{1}{d}\sum_{i=1}^{d}x_i,\\
 & \sigma=\sqrt{\tfrac{1}{d}\sum_{i=1}^{d}(x_i-\mu)^2}
\end{aligned}$\\
 & RMSNorm [5] &
$\tfrac{x}{\operatorname{RMS}(\symbf{x})}\cdot\gamma,\quad
 \operatorname{RMS}(\symbf{x})=\sqrt{\tfrac{1}{d}\sum_{i=1}^{d}x_i^2}$\\
 & DeepNorm [6] & $\operatorname{LayerNorm}(\alpha\cdot\symbf{x}+\operatorname{Sublayer}(\symbf{x}))$\\
\midrule
激活函数 & ReLU [7] & $\operatorname{ReLU}(\symbf{x})=\max(0,\symbf{x})$\\
 & GeLU [8] &
$\begin{aligned}[t]
\operatorname{GeLU}(\symbf{x})={}&0.5\symbf{x}\otimes{}\\
 & \left(1+\tanh\left(\sqrt{\tfrac{2}{\pi}}\,(x+0.044715x^3)\right)\right)
\end{aligned}$\\
 & Swish [9] & $f(x)=x\cdot\tfrac{1}{1+e^{-x}}$\\
 & SwiGLU [10] & $f(x)=x\odot\sigma(Wx+b)$\\
 & GeGLU [10] & 与 SwiGLU 类似，采用 GeLU\\
\midrule
位置嵌入 & 绝对位置 [1] & $x_i=x_i+p_i$\\
 & 相对位置 [11] & $A_{ij}=W_qx_ix_j^{\mathrm{T}}W_k+r_{i-j}$\\
 & RoPE [12] & $A_{ij}=W_qx_iR_{\theta,i-j}x_j^{\mathrm{T}}W_k$\\
 & Alibi [13] & $A_{ij}=W_qx_ix_j^{\mathrm{T}}W_k-m(i-j)$\\
\bottomrule
\end{tabularx}

\medskip
\begin{minipage}{\linewidth}
[1] Vaswani et al. (2017)，[2] Radford et al. (2019)，[3] Ding et al. (2021)，[4] Ba et al. (2016)，[5] Zhang and Sennrich (2019)，[6] Wang et al. (2022)，[7] Nair and Hinton (2010)，[8] Wang et al. (2019)，[9] Ramachandran et al. (2017)，[10] Shazeer (2020)，[11] Raffel et al. (2020)，[12] Su et al. (2021)，[13] Press et al. (2021)。
\par\smallskip
\textbf{译注}\quad LayerNorm 一行按原书印刷形式使用 $\sqrt{\sigma}$ 作分母，并用 $\sigma$ 表示标准差。正文式（2.18）的分母为 $\sigma$。
\end{minipage}
\end{table}


\paragraph{结构改动}
为了应对 Transformer 的计算需求，研究者提出多种整体结构上的改动。Lite Transformer 引入双分支结构，用注意力处理长程上下文，用卷积处理局部依赖，适用于移动设备。Funnel Transformer 引入漏斗状的编码器架构，DeLighT 则以专门模块取代标准 Transformer 块，分别力图降低浮点运算量（FLOPs）、内存占用和模型大小。由于计算复杂度呈二次增长，Transformer 难以处理长序列。循环式和分层式 Transformer 等分治策略应运而生。循环式 Transformer（recurrent Transformer）用缓存存储历史信息，Transformer-XL 等技术由此延长上下文。分层式 Transformer（hierarchical Transformer）把输入分成更小的片段，先处理低层特征，再聚合起来作高层处理，有助于处理长输入并生成更丰富的表示。

% PDF 73--78；原书第 43--48 页。
\section{数据}
\label{sec:ch02-data}
前面主要讨论了大语言模型的技术概念。除了模型架构本身，要理解大语言模型如何工作，训练数据也同样重要。本节介绍常见的训练数据类型及其对大语言模型能力的影响。

\subsection{语言模型的预训练数据集}
自 2018 年以来，迁移学习成为自然语言处理各领域的主流。那一年，先后发布了 ULMFiT、GPT 和 BERT 三个重要的语言模型。三者的架构差异很大，但有一项共同做法，即只使用自监督语言建模目标（self-supervised language modeling objective）进行预训练，再用特定任务的有标签数据微调。由此可以借助大量文本获得一般的语言理解能力，无需为文本标注，而有标签数据往往难以取得。本节介绍语言模型预训练最常用的数据来源。

预训练的目标是让大语言模型具备通用的语言理解能力和世界知识。因此，选用的训练数据应覆盖广泛的主题，使用丰富的词汇，并包含具有代表性的书面语言模式分布。数据当然还必须足够多。常用来源如下。
\begin{itemize}
\item \textbf{网页抓取（web scraping）。}从某个网页中的链接出发自动收集网页，再沿新收集网页中的链接继续收集，依此类推。这类数据语言范围广，但质量参差不齐。互联网上的俚语、拼写错误和其他不规范表达，有助于模型适应不同输入。另一方面，许多文本可能无法理解或与事实相悖，若清理不充分，会给模型带来不利影响。Common Crawl 是最知名的公开网页抓取数据。
\item \textbf{维基百科。}维基百科数据有多项好处。首先，它提供大量事实信息，通常编辑得较好，格式也统一，较少出现普通网页中的数据质量问题。维基百科还收录多种语言的文章，便于将训练范围扩展到英语以外。
\item \textbf{图书。}小说提供了优质的叙事文本，描述人们如何思考，如何与他人及环境互动。这类语言在维基百科等知识库中找不到，后者只用第三人称叙述事件。大多数图书也很适合用来建模长程依赖。显然，故事书中的许多信息是虚构的。
\item \textbf{代码。}随着生成模型能力增强，代码生成成为常见应用。人们经常利用 GitHub 和 StackExchange 的数据训练能生成代码的模型。值得注意的是，使用代码训练也可能增强大语言模型在其他逻辑推理任务中的能力（\refcite{ch02-fu}{Fu and Khot, 2022}）。
\end{itemize}

按早期 Transformer 模型的训练规模，通常选用上述来源中的一两种。到了现代大语言模型的规模，更常见的做法是把这些来源乃至更多来源合在一起，利用各自的长处。The Pile（\refcite{ch02-gao}{Gao et al., 2020}）汇集了 22 个来源的语料，包括法律文本、医学文本、学术论文以及 GitHub 上的代码。研究表明，与在 Common Crawl 等来源较单一的语料上训练的模型相比，使用这些数据可以改善下游任务表现。ROOTS 语料库（\refcite{ch02-laurencon}{Laurençon et al., 2023}）把这一思路进一步扩展，从数百个来源纳入 46 种自然语言和 13 种编程语言。

\begin{table}[!htbp]
\centering
\caption{大语言模型预训练中广泛采用的若干语料库。}
\label{tab:ch02-corpora}
\small
\begin{tabularx}{\linewidth}{@{}lX@{}}
\toprule
语料库 & 来源\\
\midrule
BookCorpus & 图书\\
Wikitext103 & 维基百科\\
Common Crawl & 互联网\\
OpenWebText & 互联网\\
The Pile & 互联网、学术研究、图书、对话、代码\\
ROOTS & 高资源和低资源语言、互联网、代码\\
\bottomrule
\end{tabularx}
\end{table}

\subsubsection{多语言语料库和平行语料库}
很多大语言模型完全或主要用一种语言训练。要让模型理解多种语言并在其间翻译，需要涵盖所需各种语言的数据。这类数据大致分为两类。
\begin{itemize}
\item \textbf{平行语料库（parallel corpus）。}每个文本样例都有一个对应的第二种语言的译文。训练使用这些语言对，把一种语言作为输入，另一种作为目标。模型预测与目标越接近，得分越高。
\item \textbf{多语言语料库（multilingual corpus）。}包含多种语言的数据，各语言之间没有明确的翻译对应关系。此类语料适用于语言建模目标，无法直接用于上述平行语料库所用的机器翻译目标。
\end{itemize}
近年来，现代大语言模型的规模已经足以让模型在少样本设置下表现出良好的翻译能力，无需特意用平行数据训练（\refcite{ch02-workshop}{Workshop et al., 2023}）。翻译能力来自模型对多种语言的共同学习，而无须从明确成对的不同语言文本中学习。

\subsection{数据预处理}
预训练所用的语料库远超人工可审阅的规模。因此，人们采用多种方法过滤可能妨碍模型表现或造成非预期影响的数据。过于偏离语言分布的文本，以及含有冒犯性内容或敏感个人信息的文本，都应剔除。

\begin{figure}[!htbp]
\centering
% 重绘自原书第 45 页的图 2.7。
\begin{tikzpicture}[
  font=\footnotesize\sffamily,
  sheet/.style={draw=black!37,fill=black!9,minimum width=14mm,minimum height=11mm,align=center},
  filterstage/.style={single arrow,draw=black!38,fill=yellow!27,minimum width=17mm,
    minimum height=8mm,single arrow head extend=1mm,inner sep=1pt,align=center}]
\node[sheet] at (.17,.18) {};
\node[sheet] at (.08,.09) {};
\node[sheet] at (0,0) {原始\\数据};
\node[filterstage] at (1.90,0) {质量过滤};
\node[filterstage] at (3.86,0) {去重};
\node[filterstage] at (5.82,0) {移除\\有害数据};
\node[filterstage] at (7.78,0) {词元切分};
\node[draw=green!42!gray,rounded corners=2.6mm,fill=green!21,
  minimum width=19mm,minimum height=9.5mm,align=center] at (9.74,0) {训练\\大语言模型};
\end{tikzpicture}

\caption{将大规模语料用于大语言模型预训练前，通常要经过的一系列步骤。}
\label{fig:ch02-cleaning}
\end{figure}

\subsubsection{低质量数据}
如图~\ref{fig:ch02-cleaning} 所示，预处理第一阶段关注数据的整体质量。原始语料通常很大，因此可以删掉相当一部分出现不适合训练迹象的数据。这一阶段的预处理可以较为粗略。

大规模语料库中常见的一个质量问题，是出现超出模型预期用途的语言。例如，若模型专门为西班牙语应用训练，出现西班牙语以外的语言会降低训练效率。可以用语言分类模型或基于规则的方法过滤这些数据。

另一项有用的预处理操作，是依据异常文本模式作统计过滤（statistical filtering）。例如，文本频繁出现远长于普通单词的字符串，标点密度很高，或者过长和过短的句子占比过高。这些模式都表明，相应文档或文档集合较难推广到其他情境，训练价值因而较低。

\subsubsection{重复数据}
人们对训练数据重复的影响已有大量讨论。\refcite{ch02-hernandez}{Hernandez 等人（2022）}发现，使用重复数据训练可能产生几种不良后果。作为对照，\refcite{ch02-biderman}{Biderman 等人（2023）}的分析表明，在重复数据上训练对模型既无益也无害。即使在最理想的情形下，用重复数据训练也似乎无法充分利用计算资源。因此，预处理时尽可能去掉重复文本，已成为标准做法。

\subsubsection{有害数据}
上述问题主要涉及只用最适合的数据来提高训练效率。还有一些信息可能不适合让模型习得。例如，如果大语言模型生成的内容中出现从网上抓取的真实姓名和电子邮件地址，就可能带来问题。训练数据中的毒性和偏见也值得关注。应对这些因素更为复杂，后续章节会再讨论，但在预处理阶段就尽可能剔除冒犯性语言，是有价值的。

\subsubsection{文本规范化}
有些数据的问题较轻，需要清理，但无须完全删去文本。例如，从网页抓取的数据自然会有残留的 HTML 标签，应将它们去掉。另一个常用步骤是 Unicode 规范化（Unicode normalization）。等价的字符串可能有多种编码方式，通常可以选用一种方法尽量统一表示，无须让模型自行学习这些等价关系。类似地，也可以按需把所有文本转为小写，使模型不再把大写字母当作不同的字符。

\subsubsection{词元切分}
预处理完成后，数据用于训练第 2.3.3 节所述的一类词元切分器（tokenizer）。这一步必须先于大语言模型训练，因为切分结果是模型的输入。常见做法是使用现有的词元切分器，免去从头训练，但前提是所用数据来源相近。首先，切分方式必须反映训练数据包含的语言，包括编程语言。此外，对话数据可能偏向“thx”“omg”一类缩写词元，而学术文献中表示专业术语的词元，其分布可能很不一样。

数据按文本块送入词元切分器，每块映射为一串词元。为提高效率，词元表示成长度为 $l$ 的整数向量，$l$ 等于子词的数量。模型第一层又称嵌入层（embedding layer），形状为 $n\times m$，其中 $n$ 是词元切分器习得的词元总数，$m$ 是预先设定的嵌入维数。因此，切分器的输出是一列查找索引，按索引为原始输入识别出的每个词元查得一个 $m$ 维向量。文本由此转换为 $l\times m$ 的浮点数矩阵，可以送入模型开始学习。

\subsection{数据对大语言模型的影响}
如前所述，训练大语言模型有多种数据来源。Gopher 的实验结果（\refcite{ch02-rae-2022}{Rae et al., 2022}）表明，调整各来源所占的数据比例，会明显影响大语言模型在多种下游任务上的整体表现。一般来说，数据多样性往往能改善许多任务的表现，同时也要考虑模型的预期用途。构建聊天机器人时，通常希望训练数据中有相当一部分是对话。反过来说，除非聊天机器人用于提供法律建议，否则没有必要纳入大量法律文档。

\begin{keypoint}
模型在预训练中接触的数据量影响很大。Chinchilla 发布时，这一点已十分明确（\refcite{ch02-hoffmann}{Hoffmann et al., 2022}）。它表明此前的大语言模型训练不足。在追求扩大模型规模所带来的能力时，人们误判了数据规模的作用。Chinchilla 研究者通过实验，尝试建立一个公式，给定算力预算后确定最合适的参数数量和训练词元数量。他们发现，模型和数据的规模应大致按比例增加，与过去强调增加参数的做法明显不同。这说明，有很强能力的大语言模型可以比先前所认为的更小。遵循这些原则，可以更有效地利用预训练预算，微调和推理的成本也会下降。
\end{keypoint}

\subsection{特定任务的数据集}
出于研究需要，自然语言处理“任务”常被用作一种通用尺度，近似衡量给定模型在各种真实场景中的表现。大部分特定任务的数据集都经过精心整理和标注，以便用于监督训练和评估。因此，它们往往比大语言模型预训练所用的超大无标签数据集小得多。

特定任务的数据集通常预先分成训练集（training set）和测试集（test set），确保所有研究者在同样的样例上训练和测试。在这种标准化的数据集上评估表现，可以直接比较不同的架构与训练策略。值得注意的是，大语言模型通常能在未见过相应训练数据的样例时，就在测试集上取得不错的评估结果。这称为零样本学习。

% PDF 78--91；原书第 48--61 页。
\section{预训练大语言模型的设计选择}
本节考察不同大语言模型在设计上的区别（\refcite{ch02-zhao}{Zhao et al., 2023}），讨论各种预训练方法和目标、Transformer 架构的选择，以及不同的解码策略。

\subsection{预训练方法}
了解各种预训练方法，对于在不同领域有效部署语言模型很重要。每种方法各有长处和难题，适用于特定任务和数据类型。下面介绍五种主要方法，说明它们如何工作、用在什么场合，以及各自的优缺点（\refcite{ch02-kalyan}{Kalyan et al., 2021}）。

\subsubsection{从头预训练}
从头预训练（Pre-training from Scratch，PTS）指用大量无标签文本，从头训练 Transformer 模型。这是构建基于 Transformer 的预训练语言模型的基础方法。这类模型通常由一个嵌入层和随后的多个 Transformer 层组成。从头预训练不依赖已有知识，因此可以作为多种应用的起点。不过，该方法需要大量算力和时间，尤其在模型和数据集都很大时。BERT、RoBERTa、ELECTRA 和 T5 等模型，都是在大量无标签文本上从头预训练的。

\subsubsection{持续预训练}
持续预训练（Continual Pre-training，CPT）是在从头预训练之后，使用特定领域的语料库进一步训练模型。此法有助于完成需要专门知识的任务，改善模型在特定领域的表现。例如，BioBERT 是 BERT 的一种变体，在生物医学文本上经过持续预训练，因而擅长生物医学和临床领域的任务。它的缺点是模型可能过拟合特定领域的语料库，从而失去部分泛化能力。

\subsubsection{同步预训练}
同步预训练（Simultaneous Pre-training，SPT）把特定领域语料与通用领域语料合在一起预训练模型。这种方法旨在取得平衡，让模型同时习得通用知识和专门知识。ClinicalBERT 是一个例子，它在临床记录与通用领域文本组成的混合语料库上预训练。同步预训练可以带来较均衡的知识结构。难点在于怎样选择和组合语料，以免模型偏向其中一个领域。

\subsubsection{任务适应性预训练}
任务适应性预训练（Task Adaptive Pre-training，TAPT）使用与任务相关的小规模语料库作预训练。它比其他方法耗费的资源少，尤其适用于特定任务可用数据有限的情形。任务适应性预训练可以补充其他方法，进一步调整经过从头预训练或持续预训练的模型，改善它们在特定任务上的表现。不过，效果在很大程度上取决于预训练语料与任务的相关程度及其质量。

\subsubsection{知识继承预训练}
知识继承预训练（Knowledge Inherited Pre-training，KIPT）是一种新方法，使用自监督学习，并从已有预训练模型中继承知识。这一思路源于人类的学习过程。人除了自学，也向有知识的人学习。知识继承预训练减少了从头预训练需要的时间和资源，因而效率较高。不过，能否成功取决于从已有模型继承的知识是否优质且相关。有效整合或迁移这些知识也未必容易。

\subsection{预训练任务}
监督学习对人工智能的发展至关重要，实际训练模型时，需要大量人工标注数据。这些模型虽然擅长特定任务，却往往需要大量有标签数据，因而训练过程费钱费时。医学、法律等专业领域中，这类数据稀缺，问题尤其明显。监督学习模型也缺乏泛化能力，往往只从给定的数据中学习，从而产生泛化误差和非预期的相关性。认识到这些局限，研究者开始探索自监督学习（Self-Supervised Learning，SSL）等其他范式。自监督学习依据数据属性和预训练任务的定义，自动生成标签。预训练任务提供伪监督信号，帮助模型学习通用知识。自监督学习主要利用大量可以自由获取的无标签数据，学习通用语言表示并提高泛化能力。

自监督学习的损失函数为
\begin{equation}
\mathcal{L}_{\mathrm{SSL}}=
\lambda_1\mathcal{L}_{\mathrm{PT}1}+
\lambda_2\mathcal{L}_{\mathrm{PT}2}+\ldots+
\lambda_m\mathcal{L}_{\mathrm{PT}m}.
\tag{2.20}
\end{equation}
其中各项含义如下。
\begin{itemize}
\item $\mathcal{L}_{\mathrm{SSL}}$ 是自监督学习的总损失函数。
\item $\mathcal{L}_{\mathrm{PT}1},\mathcal{L}_{\mathrm{PT}2},\ldots,\mathcal{L}_{\mathrm{PT}m}$ 分别是各项预训练任务的损失函数。
\item $\lambda_1,\lambda_2,\ldots,\lambda_m$ 是分配给各预训练任务损失的权重，决定各自对总损失的贡献。
\end{itemize}
为训练不同的大语言模型，已有研究提出了许多自监督预训练任务（\refcite{ch02-kalyan}{Kalyan et al., 2021}）。下面考察其中一些常见任务。

\subsubsection{因果语言模型}
因果语言建模（Causal Language Modeling，CLM）根据上下文预测序列中的下一个词，可以从左到右，也可以从右到左。给定序列 $x=\{x_1,x_2,x_3,\ldots,x_{|x|}\}$，其中 $|x|$ 表示词元数量，CLM 的损失函数定义为
\begin{equation}
\mathcal{L}_{\mathrm{CLM}}^{(x)}=-\frac{1}{|x|}\sum_{i=1}^{|x|}\log P(x_i\mid x_{<i}),
\tag{2.21}
\end{equation}
其中 $x_{<i}$ 表示序列中先于 $x_i$ 的词元。

\subsubsection{掩码语言模型}
掩码语言建模（Masked Language Modeling，MLM）用于预训练阶段。它遮盖输入序列中选中的词元，再训练模型预测这些词元。令 $x_{\backslash M_x}$ 表示 $x$ 的掩码版本，$M_x$ 表示 $x$ 中被遮盖词元的位置集合。MLM 的损失定义为
\begin{equation}
\mathcal{L}_{\mathrm{MLM}}^{(x)}=-\frac{1}{|M_x|}\sum_{i\in M_x}\log P(x_i/x_{\backslash M_x}).
\tag{2.22}
\end{equation}
模型学习准确预测被遮盖的词元，使该损失最小，从而加深对语言结构的理解。BERT 是自然语言处理中的重要模型，它采用 MLM 作为预训练任务，以 0.15 的概率选取要遮盖的词元。

\subsubsection{替换词元检测}
替换词元检测（Replaced Token Detection，RTD）通过增强训练信号、缩小预训练与微调之间的差异，缓解 MLM 的局限。MLM 使用特殊掩码词元来破坏输入，而 RTD 用以 MLM 目标预训练的模型生成的词元来破坏句子。任务由此变为词元级二分类，判断每个词元是否被替换。过程分为两步，先用 MLM 目标训练生成器模型，再用 RTD 目标训练以生成器初始化的判别器模型。RTD 的损失写为
\begin{equation}
\mathcal{L}_{\mathrm{RTD}}^{(x)}
=-\frac{1}{|\hat{x}|}\sum_{i=1}^{|\hat{x}|}\log P(d/\hat{x}_i),
\tag{2.23}
\end{equation}
其中 $d\in\{0,1\}$ 表示词元是否被替换，替换为 1，未替换为 0。$\hat{x}$ 是被破坏的句子，$P(d/\hat{x}_i)$ 表示词元被替换或未被替换的概率。

\subsubsection{打乱词元检测}
打乱词元检测（Shuffled Token Detection，STD）旨在提高模型理解连贯句子结构的能力，进而改善模型在多种任务上的表现。该任务以 0.15 的概率打乱序列中的词元，其损失为
\begin{equation}
\mathcal{L}_{\mathrm{STD}}^{(x)}
=-\frac{1}{|\hat{x}|}\sum_{i=1}^{|\hat{x}|}\log P(d/\hat{x}_i).
\tag{2.24}
\end{equation}
式中 $d\in\{0,1\}$ 表示词元是否被替换，替换为 1，未替换为 0，$\hat{x}$ 是被破坏的句子。模型通过学习在序列上下文中有效辨识和理解被打乱的词元，使这一损失最小。
\transnote{原书在 STD 的说明中仍用“替换”描述 $d$ 的取值。}

\subsubsection{随机词元替换}
随机词元替换（Random Token Substitution，RTS）由 \refcite{ch02-liello}{Liello 等人（2021）}提出，用来识别序列中被随机换掉的词元。该技术从词表中选择其他词元，随机替换给定序列中 15\% 的词元。这个方法效率较高，因为无须另设生成器模型破坏输入序列。RTS 的损失为
\begin{equation}
\mathcal{L}_{\mathrm{RTS}}^{(x)}
=-\frac{1}{|\hat{x}|}\sum_{i=1}^{|\hat{x}|}\log P(d/\hat{x}_i),
\tag{2.25}
\end{equation}
其中 $d\in\{0,1\}$ 表示词元是否被随机换掉，换掉为 1，未换掉为 0。$\hat{x}$ 是将原序列 $x$ 中 15\% 的词元随机换掉后所得的序列。

\subsubsection{替换式语言建模}
替换式语言建模（Swapped Language Modeling，SLM）处理 MLM 预训练任务中使用特殊掩码词元所造成的差异。这个差异出现在预训练与微调两个阶段之间。SLM 以 0.15 的概率用从词表随机选取的词元替换输入，从而减轻这种差异。SLM 与 MLM 都预测被破坏的词元，但 SLM 用随机词元替换原词元，MLM 则用掩码词元。SLM 和 RTS 都用随机词元破坏输入，但 SLM 的样本效率较低。这是因为 SLM 只在 15\% 的输入词元上计算目标，RTS 则使用序列中的每个词元。SLM 的损失定义为
\begin{equation}
\mathcal{L}_{\mathrm{SLM}}^{(x)}
=-\frac{1}{|R_x|}\sum_{i\in R_x}\log P(x_i/x_{\backslash R_x}),
\tag{2.26}
\end{equation}
其中 $R_x$ 是被随机替换词元的位置集合，$x_{\backslash R_x}$ 是 $x$ 的破坏版本。

\subsubsection{翻译语言建模}
翻译语言建模（Translation Language Modeling，TLM）用于多语言模型的预训练。给定两个语言不同的句子，TLM 遮盖两个句子中的一部分词元，再训练模型预测这些词元。TLM 的损失定义为
\begin{equation}
\begin{aligned}
\mathcal{L}_{\mathrm{TLM}}^{(x)}
={}&-\frac{1}{|M_x|}\sum_{i\in M_x}\log P(x_i/x_{\backslash M_x},y_{\backslash M_y})\\
 &-\frac{1}{|M_y|}\sum_{i\in M_y}\log P(y_i/x_{\backslash M_x},y_{\backslash M_y}).
\end{aligned}
\tag{2.27}
\end{equation}
其中 $M_x$ 和 $M_y$ 分别表示句子 $x$ 和 $y$ 中被遮盖位置的集合，$x_{\backslash M_x}$ 和 $y_{\backslash M_y}$ 分别是 $x$ 和 $y$ 的掩码版本。

\subsubsection{交替语言建模}
交替语言建模（Alternate Language Modeling，ALM）用于跨语言模型预训练。其做法是改变输入序列中各句子使用的语言。给定一对平行句子 $(x,y)$，随机选取 $x$ 中的一些短语，用 $y$ 中对应的译文替换，生成一个语码转换（code-switched）的句子。ALM 遵循标准 MLM 的掩码流程，选取要遮盖的词元。在这些语码转换句子上预训练模型，可以更有效地学习语言之间的关系。
\begin{equation}
\mathcal{L}_{\mathrm{ALM}}^{(z(x,y))}
=-\frac{1}{|M|}\sum_{i\in M}\log P(z_i/z_{\backslash M}).
\tag{2.28}
\end{equation}
其中 $z$ 是根据 $x$ 和 $y$ 生成的语码转换句子，$z_{\backslash M}$ 是 $z$ 的掩码版本，$M$ 是 $z_{\backslash M}$ 中被遮盖词元的位置集合。

\subsubsection{句段边界目标}
句段边界目标（Sentence Boundary Objective，SBO）根据片段边界词元和位置嵌入预测被遮盖的词元。SBO 的损失定义为
\begin{equation}
\mathcal{L}_{\mathrm{SBO}}^{(x)}
=-\frac{1}{|S|}\sum_{i\in S}\log P(x_i/f(x_{s-1},x_{e+1},p_{s-e+1})),
\tag{2.29}
\end{equation}
其中 $f(\cdot)$ 是两层前馈神经网络，$S$ 表示连续片段内的词元位置，$s$ 和 $e$ 分别表示片段的起点和终点，$p$ 表示位置嵌入。

\subsubsection{下一句预测}
下一句预测（Next Sentence Prediction，NSP）是对句子对作二分类的任务，其损失函数定义为
\begin{equation}
\mathcal{L}_{\mathrm{NSP}}^{(x,y)}=-\log P(d/x,y),
\tag{2.30}
\end{equation}
其中 $d$ 是二值变量，表示句子 $(x,y)$ 是否连续。连续为 1，不连续为 0。

\subsubsection{句序预测}
句序预测（Sentence Order Prediction，SOP）关注句子是否连贯。下一句预测与它不同，还包含主题预测。ALBERT 引入了句序预测，用来判断句子是否按正确顺序排列，还是已经调换。训练样例保持平衡，有 50\% 的样例调换了顺序。SOP 的损失为
\begin{equation}
\mathcal{L}_{\mathrm{SOP}}^{(x,y)}=-\log P(d/x,y),
\tag{2.31}
\end{equation}
其中 $d\in\{1,0\}$ 表示句子是否调换了顺序。

\subsubsection{序列到序列语言建模}
序列到序列语言建模（Sequence-to-Sequence Language Modeling，Seq2Seq）是 MLM 的扩展，用于预训练基于编码器—解码器的模型。损失定义为
\begin{equation}
\mathcal{L}_{\mathrm{Seq2Seq}}^{(x)}
=-\frac{1}{l_s}\sum_{s=i}^{j}\log P(x_s/\hat{x},x_{i:s-1}),
\tag{2.32}
\end{equation}
其中 $\hat{x}$ 是 $x$ 的掩码版本，$l_s$ 表示被遮盖的 $n$ 元片段的长度。

\subsubsection{去噪自编码器}
去噪自编码器（Denoising Autoencoder，DAE）从被破坏的文本重建原始文本。DAE 的损失定义为
\begin{equation}
\mathcal{L}_{\mathrm{DAE}}
=-\frac{1}{|x|}\sum_{i=1}^{|x|}\log P(x_i/\hat{x},x_{<i}),
\tag{2.33}
\end{equation}
其中 $\hat{x}$ 是 $x$ 的破坏版本。

% PDF 85--91；原书第 55--61 页。
\subsection{架构}
\refcite{ch02-vaswani}{Vaswani 等人（2017）}最初提出的 Transformer，由堆叠的编码器层和解码器层组成。预训练基于 Transformer 的语言模型时，可以堆叠编码器层、解码器层，也可以两者并用，由此得到多种架构，如图~\ref{fig:ch02-patterns} 所示。

\subsubsection{编码器—解码器}
编码器—解码器架构分为两部分。编码器处理输入序列，将输入变成一个包含所学输入信息的连续表示。解码器再用这个表示生成输出序列。该架构适用于机器翻译、文本摘要等序列到序列任务。例如，在机器翻译中，编码器处理源语言句子，解码器生成目标语言译文。该架构的注意力机制使模型在生成输出时能够关注输入序列的不同部分，动态计算所需的上下文。

\subsubsection{因果解码器}
因果解码器架构面向自回归任务，模型在这类任务中逐个生成词元。它采用单向注意力机制，即生成时每个词元只能关注前面的词元及其自身。对于需要生成连贯且符合上下文的文本生成任务，这一点尤其有用。例如，补全文本时，模型根据前面的词元预测下一个词元，使生成的文本连贯并与上下文相关。

\begin{figure}[!htbp]
\centering
% 重绘自原书第 56 页的图 2.8；轴标签是原示例中的六个英文词元。
\begingroup
\colorlet{prefixcell}{blue!24}
\colorlet{crosscell}{green!32!gray!13}
\colorlet{targetcell}{yellow!37!orange!12}
\colorlet{maskedcell}{blue!12!gray!37}
\newcommand{\attentiontiles}[2]{%
  \node[font=\footnotesize\sffamily\bfseries] at (1.32,.49) {#2};
  \foreach \row/\ca/\cb/\cc/\cd/\ce/\cf in {#1} {
    \foreach \col/\tone in {0/\ca,1/\cb,2/\cc,3/\cd,4/\ce,5/\cf} {
      \node[tile,fill=\tone] at ({.53*\col},{-.53*\row}) {};
    }
  }
  \foreach \word [count=\idx] in {A,Survey,of,Large,Language,Models} {
    \node[anchor=east] at (-.24,{-.53*(\idx-1)}) {\word};
    \node[rotate=90,anchor=east] at ({.53*(\idx-1)},-2.91) {\word};
  }
}
\newcommand{\decoderbrackets}{%
  \draw[group] (-1.60,.17) -- (-1.73,.17) -- (-1.73,-2.82) -- (-1.60,-2.82);
  \node[rotate=90] at (-1.97,-1.32) {解码器};
  \draw[group] (-.19,-3.94) -- (-.19,-4.08) -- (2.83,-4.08) -- (2.83,-3.94);
  \node at (1.32,-4.29) {解码器};
}
\begin{tikzpicture}[
  font=\scriptsize\sffamily,
  tile/.style={draw=black!56,minimum size=4.1mm,inner sep=0pt,rounded corners=1pt,line width=.48pt},
  group/.style={draw=black!49,line width=.52pt,rounded corners=1pt}]
\begin{scope}[shift={(2.74,5.13)}]
\attentiontiles{
  0/prefixcell/maskedcell/maskedcell/maskedcell/maskedcell/maskedcell,
  1/prefixcell/prefixcell/maskedcell/maskedcell/maskedcell/maskedcell,
  2/prefixcell/prefixcell/prefixcell/maskedcell/maskedcell/maskedcell,
  3/crosscell/crosscell/crosscell/targetcell/maskedcell/maskedcell,
  4/crosscell/crosscell/crosscell/targetcell/targetcell/maskedcell,
  5/crosscell/crosscell/crosscell/targetcell/targetcell/targetcell%
}{因果解码器}
\decoderbrackets
\end{scope}
\begin{scope}
\attentiontiles{
  0/prefixcell/prefixcell/prefixcell/maskedcell/maskedcell/maskedcell,
  1/prefixcell/prefixcell/prefixcell/maskedcell/maskedcell/maskedcell,
  2/prefixcell/prefixcell/prefixcell/maskedcell/maskedcell/maskedcell,
  3/crosscell/crosscell/crosscell/targetcell/maskedcell/maskedcell,
  4/crosscell/crosscell/crosscell/targetcell/targetcell/maskedcell,
  5/crosscell/crosscell/crosscell/targetcell/targetcell/targetcell%
}{前缀解码器}
\decoderbrackets
\end{scope}
\begin{scope}[xshift=5.50cm]
\attentiontiles{
  0/prefixcell/prefixcell/prefixcell/maskedcell/maskedcell/maskedcell,
  1/prefixcell/prefixcell/prefixcell/maskedcell/maskedcell/maskedcell,
  2/prefixcell/prefixcell/prefixcell/maskedcell/maskedcell/maskedcell,
  3/crosscell/crosscell/crosscell/targetcell/maskedcell/maskedcell,
  4/crosscell/crosscell/crosscell/targetcell/targetcell/maskedcell,
  5/crosscell/crosscell/crosscell/targetcell/targetcell/targetcell%
}{编码器—解码器}
\draw[group] (-1.60,.17) -- (-1.73,.17) -- (-1.73,-1.26) -- (-1.60,-1.26);
\draw[group] (-1.60,-1.38) -- (-1.73,-1.38) -- (-1.73,-2.82) -- (-1.60,-2.82);
\node[rotate=90] at (-1.97,-.55) {编码器};
\node[rotate=90] at (-1.97,-2.10) {解码器};
\draw[group] (-.19,-3.94) -- (-.19,-4.08) -- (1.26,-4.08) -- (1.26,-3.94);
\draw[group] (1.38,-3.94) -- (1.38,-4.08) -- (2.83,-4.08) -- (2.83,-3.94);
\node at (.53,-4.29) {编码器};
\node at (2.12,-4.29) {解码器};
\end{scope}
\end{tikzpicture}
\endgroup

\caption{三种主要架构的注意力模式。蓝色、绿色、黄色和灰色圆角方块分别表示前缀词元内部的注意力、前缀词元与目标词元之间的注意力、目标词元内部的注意力以及被掩码遮盖的注意力。}
\label{fig:ch02-patterns}
\end{figure}

\subsubsection{前缀解码器}
前缀解码器是因果解码器的变体。模型可以在一段作为前缀的词元内部双向计算注意力，对其余词元仍保持单向注意力。采用这种混合机制，模型每生成一个词元时都可以利用更广的上下文，因而适合需要理解序列中前后词元的任务。例如，在对话系统中，生成下一个词元时，模型可以关注对话历史和已经生成的部分回复。

\subsubsection{编码器}
编码器能够高效处理并理解输入序列中的上下文信息，适合某些自然语言处理任务。各编码器层都会为输入序列生成能够反映上下文的表示。最后一个编码器层的最终输出用作上下文表示，为多种下游任务提供输入。编码器架构特别适合需要深入理解词元上下文、无须生成序列的任务，例如分类任务。

\subsubsection{混合专家}
混合专家（Mixture of Experts，MoE）架构是 Transformer 的一种变体。它用混合专家层替换标准前馈模块，如图~\ref{fig:ch02-moe} 所示。每个混合专家层包含多个并行单元，称为“专家”，每个专家都有各自的参数。路由器（router）根据专家的能力，把输入词元交给相应专家。专家是注意力模块后的前馈层，彼此独立地处理词元。传统模型扩大容量会增加计算成本，而混合专家架构每次只激活少量专家。这种稀疏激活让模型规模能够继续增大，而计算需求无需等比例增加，因而仍能高效运行。

\sourcefigure{fig02_09.pdf}{Transformer 架构的混合专家变体。左侧路由器（Router）将输入 $x_1$ 送至 FFN 2（$p=0.65$），将输入 $x_2$ 送至 FFN 1（$p=0.80$）。右侧标出混合专家层（MoE layer）、层归一化（Layer Norm）和掩码自注意力（Masked Self-Attention）的位置。}{fig:ch02-moe}

\subsection{大语言模型预训练的技巧与策略}
本节讨论训练大语言模型时常用的配置、方法和策略。

\subsubsection{训练方法}
\begin{itemize}
\item \textbf{学习率。}大多数大语言模型预训练时采用相近的学习率计划，分为预热（warm-up）和衰减阶段。先用约 0.1\% 至 0.5\% 的训练步数逐渐增大学习率，学习率通常在 $5\times 10^{-5}$ 至 $1\times 10^{-4}$ 的范围。此后采用余弦衰减（cosine decay）策略逐步减小学习率。
\item \textbf{批量大小。}语言模型预训练常用大批量，如 2\,048 个样例或 400 万个词元，以改善训练的稳定性和效率。GPT-3 和 PaLM 等模型会在训练中动态调整批量大小。例如，GPT-3 的批量大小从 3.2 万词元增至 320 万词元。研究表明，适应性调整批量大小能有效使大语言模型训练更平稳。
\item \textbf{优化器。}训练 GPT-3 等大语言模型时，常使用 Adam 和 AdamW。这些优化器根据梯度估计调整更新，典型超参数（hyperparameter）为 $\beta_1=0.9$、$\beta_2=0.95$ 和 $\epsilon=10^{-8}$。PaLM、T5 等模型还采用 Adafactor 优化器，这是内存用量较少的 Adam 变体。它的超参数 $\beta_1=0.9$，$\beta_2$ 则随训练步数调整。
\end{itemize}

\subsubsection{解码策略}
\paragraph*{贪心搜索}
自回归解码（autoregressive decoding）是纯解码器架构所用的技术之一。贪心搜索（greedy search）是最常用的解码方法之一。它在每一步根据此前生成的词元，预测概率最大的词元，数学形式如下。
\[
x_i=\arg\max_x P(x\mid x_{<i}),
\]
其中 $x_i$ 是第 $i$ 步预测的词元，即给定上下文 $x_{<i}$ 时概率最大的词元。以未完成的句子“The sky is so”为例，“blue”很可能恰当地补全这个句子，贪心搜索可能把它预测为下一个词元。在机器翻译和文本摘要等文本生成任务中，输入与预期输出之间有较强的依赖，这种办法很有效。

在输出需要紧扣输入的情形中，贪心搜索利用概率和上下文给出可靠的结果。这种解码策略也可以用于编码器—解码器和前缀解码器模型。研究者提出了许多改进，下面讨论其中一些。束搜索（beam search）是一种常用策略，在每一步解码时保留概率最大的 $n$ 个句子，最终选择其中概率最高的一个。

\begin{practicaltip}
束宽通常取 3 至 6 就已足够。继续增大束宽，表现可能反而下降。长度惩罚（length penalty），也称长度归一化（length normalization），是另一项改进，用于补偿束搜索偏好较短句子的倾向。做法是按句子长度调整句子概率，以长度的某个幂作除数。也有方法对生成已用过的词元施加惩罚，以减少重复词元或重复 $n$ 元组。多样化束搜索（diverse beam search）也是一种改进，可以对同一个输入产生多种不同的输出。
\end{practicaltip}

\paragraph*{随机搜索}
基于采样的解码策略以概率方式选择词元，为生成的文本增加多样性与随机性。如果目标是生成丰富且有吸引力的文本，这种策略很有用。例如，给定“I am thirsty. I would like a cup of”作为上下文，下一个词元的概率分布可能偏向“tea”“coffee”或“water”。基于采样的方法仍可能选到概率较低的词，只是机会较小。虽然“tea”的概率最高，“coffee”“water”和“juice”等词仍有机会被选中，使回答更多样。这种方法可用于纯解码器、编码器—解码器和前缀解码器等架构，可灵活适应不同的语言生成任务。

人们对随机采样作了多种改进，减少选中极低概率词的机会，从而改善生成文本的质量。一种改进是温度采样（temperature sampling），计算词表中各词元的概率时，调整 softmax 函数中的温度系数。
\[
P(x_j\mid x_{<i})=\frac{\exp(l_j/t)}{\sum_{j'}\exp(l_{j'}/t)}.
\]
其中 $l'_j$ 表示各词的 logit，$t$ 是温度系数。温度越低，概率较高的词越容易被选中，概率较低的词越不容易被选中。温度为 1 时，默认按原概率分布随机采样。温度趋近 0 时，结果接近贪心搜索。温度无限增大时，采样趋近均匀采样。
\transnote{原书第 59 页的公式用 $l_j$ 表示 logit，紧随其后的文字写作 $l'_j$。}

\begin{practicaltip}
另一项改进是 Top-$k$ 采样（Top-$k$ sampling），截去概率较低的词元，只在概率最高的 $k$ 个词元中采样。Top-$p$ 采样（Top-$p$ sampling），也称核采样（nucleus sampling），是另一种策略。它从累计概率大于或等于指定值 $p$ 的最小词元集合中采样。构造该集合时，按生成概率从高到低加入词元，直到累计概率超过 $p$。例如，按序加入词元，直到累计概率超过 0.8，就只在已加入的词元中采样。
\end{practicaltip}

\subsubsection{三维并行}
三维并行（3D parallelism）把数据并行、流水线并行和张量并行三种主要技术结合起来，以高效训练大语言模型。

\paragraph*{数据并行}
数据并行（data parallelism）把模型参数和训练数据集分布到多个 GPU，以加快训练。每个 GPU 处理分配给自己的数据并计算梯度，再汇总各设备的梯度，用于更新每个 GPU 上的模型。

\begin{practicaltip}
DeepSpeed 程序库引入的 ZeRO 技术解决数据并行中的内存冗余。通常，数据并行使每个 GPU 都要存储同一个大语言模型的完整副本，包括模型参数、梯度和优化器参数（\refcite{ch02-rajbhandari}{Rajbhandari et al., 2020}）。这种重复存储消耗额外内存。ZeRO 让每个 GPU 只保留一部分数据，需要时从其他 GPU 取回其余数据。根据分片存储的内容，提出了三种策略，即优化器状态分片、梯度分片和参数分片。测试表明，前两种不会增加通信开销。第三种使通信开销增加约 50\%，但节省的内存量取决于 GPU 数量。PyTorch 也引入了类似 ZeRO 的全分片数据并行（Fully Sharded Data Parallel，FSDP）技术。
\end{practicaltip}

\paragraph*{流水线并行}
流水线并行（pipeline parallelism）把大语言模型的不同层分配到若干 GPU。为减少数据传输开销，连续的层会分到同一个 GPU。简单的实现可能无法充分利用 GPU。GPipe 和 PipeDream 等进阶方法同时处理多个数据批次并异步更新梯度，以提高效率（\refcite{ch02-harlap}{Harlap et al., 2018}；\refcite{ch02-huang}{Huang et al., 2019}）。

\paragraph*{张量并行}
张量并行（tensor parallelism）将大语言模型的张量或参数矩阵分割，分配到多个 GPU。例如，可在矩阵乘法中按列拆分参数矩阵，交给不同的 GPU 计算，然后合并各 GPU 的结果。Megatron-LM 等程序库支持张量并行，可用于更复杂的张量（\refcite{ch02-shoeybi}{Shoeybi et al., 2019}）。

% PDF 91--105；原书第 61--75 页。
\section{常用的预训练大语言模型}
本节深入考察三种主要的大语言模型架构，讨论它们采用的数据集、与 Transformer 架构的关系、研究所得的主要结论，以及各自的变体。
\transnote{原书在本节开头称将介绍三种架构，下文还列有第 2.6.4 节的混合专家模型 Mixtral。}

\subsection{BERT（编码器）}
\label{sec:ch02-bert}
BERT，即 Transformer 双向编码器表示（Bidirectional Encoder Representation from Transformer），是一种使用注意力机制更好理解语言上下文的预训练模型（\refcite{ch02-devlin}{Devlin et al., 2019}）。BERT 由多个编码器部分组成，每一部分都为模型的稳健性作出贡献。BERT 问世时，在多种自然语言处理任务上刷新了基准成绩，包括 SQuAD v1.1 数据集上的问答和 MNLI 数据集上的自然语言推断（Natural Language Inference，NLI）。传统语言模型沿单一方向处理文本序列，而 BERT 的双向训练能让模型更全面地理解语言上下文和序列的衔接。

\subsubsection{数据集}
BERT 的训练数据主要来自维基百科，约有 25 亿词，以及 BooksCorpus，约有 8 亿词。

\subsubsection{架构}
BERT 是一个纯编码器 Transformer。它有不同架构规模的预训练模型，下面是两个例子。
\begin{itemize}
\item BERT-BASE 由 12 层组成，含 768 个隐藏单元、12 个注意力头和 1.10 亿个参数。
\item BERT-LARGE 是较大的版本，含 24 层、1\,024 个隐藏单元、16 个注意力头和 3.40 亿个参数。
\end{itemize}
BERT-BASE 使用 4 个云端张量处理单元（Tensor Processing Unit，TPU）训练了 4 天。BERT-LARGE 使用 16 个 TPU，训练时间相同。

\subsubsection{训练}
如图~\ref{fig:ch02-bert} 所示，BERT 的训练分为预训练和微调两个阶段。预训练先在各种任务中用无标签数据学习，微调则以预训练取得的参数为起点，使用目标任务的有标签数据进一步优化参数。

BERT 的训练方法结合了两个目标，即掩码语言建模和下一句预测。训练时最小化这两个目标合成的损失函数。对 BERT 来说，每个训练样例是一对句子。两个句子在原文中可能连续，也可能不连续。特殊词元 \texttt{[CLS]} 表示序列起始，\texttt{[SEP]} 表示句子之间的分隔。训练样例中的一部分词元会被 \texttt{[MASK]} 遮盖，或被随机词元替换。送入 BERT 前，词元先转为嵌入向量，再加上位置编码。BERT 还为它们加上分段嵌入（segment embedding），标明词元属于第一个还是第二个句子。

预训练完成后，BERT 可以适配多种下游任务，包括输入为单段文本或一对文本的任务。BERT 在 2\,500 亿个词元上训练 3.50 亿个参数，由此获得的通用语言表示，显著提升了多种自然语言处理任务的先进水平。微调时，可以给 BERT 加入额外的层，并使用特定下游任务的训练数据，更新新增层和预训练的 BERT 参数。Transformer 编码器实质上是预训练 BERT，接收一个文本序列，使用 \texttt{[CLS]} 的表示来预测。例如，在情感分析或分类任务中，\texttt{[CLS]} 被替换为实际的分类标签。微调时，采用基于梯度的方法最小化预测与真实标签之间的交叉熵损失。新增层从头训练，预训练的 BERT 参数也随之更新。

\begin{figure}[!tbp]
\centering
\includegraphics[width=\linewidth,height=.80\textheight,keepaspectratio]{assets/fig02_10.pdf}
\caption{BERT 可将预训练目标适配为使用特定任务输入数据的微调目标。上半部分是无标签句子对的预训练，NSP 表示下一句预测，Mask LM 表示掩码语言建模。下半部分为微调，SQuAD 示例以问题和段落为输入，预测答案片段的起点和终点，后方还示意 NER 与 MNLI 任务。}
\label{fig:ch02-bert}
\end{figure}

\subsubsection{要点}
\begin{keypoint}
\begin{enumerate}
\item 模型规模很重要。BERT-LARGE 含 3.45 亿个参数，是这一类模型中最大的。虽然与 BERT-BASE 的结构相同，它的表现更好，而 BERT-BASE “只有”1.10 亿个参数。
\item 训练数据足够时，增加训练步数与准确率的提高有关。例如，在 MNLI 任务上，批量大小均为 128\,000 词，BERT-BASE 从训练 50 万步改为训练 100 万步，准确率提高了 1.0\%。
\item 由于 BERT 的双向方法（MLM）每个批次只预测 15\% 的词，收敛可能比单向的从左到右方法慢，但经过有限次数的预训练迭代后，表现会超过后者。
\end{enumerate}
\end{keypoint}
\transnote{原书在第 62 页的架构列表、训练说明和第 64 页的要点中分别给出 3.40 亿、3.50 亿和 3.45 亿个参数，译文按各处印刷值保留。}

\subsubsection{变体}
BERT 后续的发展和变体，旨在改进模型架构或预训练目标（\refcite{ch02-kamath}{Kamath et al., 2022}）。主要例子包括以下几种。
\begin{itemize}
\item \textbf{RoBERTa} 是同等规模的 BERT 变体，在 2\,000 亿个词元上预训练。研究者发现，在这种情况下，BERT 所用的损失函数影响较小。
\item \textbf{ALBERT} 通过要求共享参数，提高效率。
\item \textbf{SpanBERT} 着重表示和预测文本片段。
\item \textbf{DistilBERT} 是利用知识蒸馏（knowledge distillation）得到的轻量版。
\item \textbf{ELECTRA} 着重进行替换词元检测。
\end{itemize}

% PDF 94--97；原书第 64--67 页。
\subsection{T5（编码器—解码器）}
文本到文本迁移 Transformer（Text-to-Text Transfer Transformer，T5）提出了一个统一框架，整合自然语言处理迁移学习过程中的多种要素（\refcite{ch02-raffel}{Raffel et al., 2020}）。其中包括不同的无标签数据集、预训练目标、基准以及微调方法。该框架用消融实验（ablation experiment）组合和比较这些要素，从中找出能够改善表现的做法。

\subsubsection{数据集}
T5 的数据来自从 Common Crawl 网页档案抽取的文本。研究者对数据作了基本的启发式过滤与预处理。提取文本后，他们剔除了不当语言、占位文本（例如 Lorem Ipsum）、代码括号（例如“\{”）、重复内容，以及没有句末标点的句子。由于主要任务面向英文文本，他们使用 \textit{langdetect7} 排除被判定为英语的置信度低于 99\% 的网页。

\subsubsection{架构}
T5 主要采用编码器—解码器架构。除少数改动外，它与原始 Transformer 的设计非常相似。T5 的一个特点是使用相对位置嵌入。它根据自注意力中键与查询之间的偏移生成可学习的嵌入，而不用固定的位置嵌入。研究提出了五种模型变体。
\begin{itemize}
\item \textbf{Base} 是与 \texttt{BERT\_base} 对应的基准模型，含 2.22 亿个参数。
\item \textbf{Small} 是 Base 的缩小版，含 6\,000 万个参数，编码器和解码器各 6 层。
\item \textbf{Large} 是 Base 的扩展版，含 7.70 亿个参数。
\item \textbf{3B} 在 Base 的基础上扩大到 30 亿个参数。
\item \textbf{11B} 是最大的变体，把 Base 扩大到 110 亿个参数。
\end{itemize}

\subsubsection{训练}
T5 采用多任务学习方法，在预训练阶段组合不同任务。按训练方式，这些任务可分为两类。
\begin{enumerate}
\item \textbf{无监督训练}
  \begin{itemize}
  \item 在 C4 数据集上用传统的语言模型训练任务和最大似然目标训练。
  \item 在掩码语言建模等无监督任务中，T5 使用从 \texttt{<extra\_id\_0>} 到 \texttt{<extra\_id\_99>} 的 100 个不同词元来组织输入与输出文本。例如，要遮盖“My name is John Smith”中的“name is”，输入变为“My \texttt{<extra\_id\_0>} John Smith”，期望输出为“\texttt{<extra\_id\_0>} name is \texttt{<extra\_id\_1>}”。
  \end{itemize}
\item \textbf{监督训练}
  \begin{itemize}
  \item 纳入问答、摘要和分类等自然语言处理任务，用整理后的数据作监督训练。所有任务都改成适合编码器—解码器模型的文本输入、文本输出形式，见图~\ref{fig:ch02-t5}.
  \item 研究采用提示方法，要求语言模型以文本给出回答。T5 处理的每项自然语言处理任务，都转换成文本到文本的格式。例如，输入可以是“\texttt{translate English to German: The house is wonderful.</s>}”，相应输出为“\texttt{<pad> Das Haus ist wunderbar.</s>}”。
  \item 研究者进行了一系列消融实验，以确定各部件的合适策略，见图~\ref{fig:ch02-t5-ablation}。最初比较三种主要方法，即（1）语言建模，（2）BERT 式掩码，（3）恢复被打乱的顺序。BERT 式方法效果最好，因此被选作后续分析。
  \item 后续实验考察了不同的输入破坏策略，如仅遮盖词元而不调换、遮盖词元并用哨兵词元（sentinel token）替代，以及删除词元。
  \end{itemize}
\end{enumerate}

\sourcefigure{fig02_11.pdf}{T5 架构把不同的自然语言处理任务统一为序列到序列任务，并按提示生成相应的回答。图中由上到下依次是德语翻译、句子可接受性判断、语义相似度评分和摘要。}{fig:ch02-t5}

\begin{figure}[!htbp]
\centering
% 重绘自原书第 67 页的图 2.12。
\begin{tikzpicture}[
  font=\footnotesize\sffamily,
  heading/.style={align=center,font=\footnotesize\sffamily\bfseries},
  choice/.style={draw=black!38,rounded corners=1.7pt,minimum height=6mm,minimum width=22mm,align=center},
  branch/.style={-{Stealth[length=1.35mm]},draw=black!58,line width=.47pt}]
\node[heading] at (0,1.47) {总体方法};
\node[heading] at (2.90,1.47) {输入破坏策略};
\node[heading] at (5.80,1.47) {破坏比例};
\node[heading] at (8.70,1.47) {受破坏的片段长度};

\node[choice,fill=blue!18] at (0,.70) {语言建模};
\node[choice,fill=blue!18] (bert) at (0,-.12) {BERT 式掩码};
\node[choice,fill=blue!18] at (0,-.94) {恢复顺序};

\node[choice,fill=yellow!27] (mask) at (2.90,.70) {遮盖};
\node[choice,fill=yellow!27] (span) at (2.90,-.12) {替换片段};
\node[choice,fill=yellow!27] (drop) at (2.90,-.94) {删除};

\foreach \idx/\y/\rate in {1/.70/10,2/-.12/15,3/-.94/25,4/-1.76/50} {
  \node[choice,fill=green!17] (rate\idx) at (5.80,\y) {\rate\%};
  \draw[branch] (span.east) to[out=0,in=180] (rate\idx.west);
}
\foreach \idx/\y/\len in {1/.70/2,2/-.12/3,3/-.94/5,4/-1.76/10} {
  \node[choice,fill=orange!22] (length\idx) at (8.70,\y) {\len};
  \draw[branch] (rate2.east) to[out=0,in=180] (length\idx.west);
}
\foreach \target in {mask,span,drop} {
  \draw[branch] (bert.east) to[out=0,in=180] (\target.west);
}
\end{tikzpicture}

\caption{用于选择 T5 训练策略的消融实验设置。}
\label{fig:ch02-t5-ablation}
\end{figure}

\subsubsection{要点}
\begin{keypoint}
\begin{itemize}
\item T5，尤其是 11B 变体，在多数自然语言处理任务中达到当时的先进水平，在 24 项任务中有 18 项领先。
\item 实验表明，为模型提供双向上下文有助于提高预测能力。
\item 破坏词的训练目标最为有效，尤其是目标序列较短的做法。这是因为用较短序列预训练，计算成本更低。
\item 过滤训练数据，尤其是移除非英语内容，已证明有益。此外，在新闻报道等特定领域数据上预训练，显著改善了相关下游任务的表现。
\item 同时用多个任务训练一个模型看似有吸引力，却导致表现下降，在 GLUE、SQuAD 和 SuperGLUE 等任务上尤其明显。
\item 加深和加宽模型、延长训练时间，可以显著改善表现。训练多个模型再使用集成方法（ensemble method），还能进一步提升表现。
\end{itemize}
\end{keypoint}

\subsubsection{变体}
\begin{enumerate}
\item \textbf{T5v1.1} 是修改架构后的增强版，只在 C4 上预训练，未加入监督任务。
\item \textbf{mT5} 是多语言变体，在涵盖 101 种语言的 mC4 语料库上训练。
\item \textbf{byT5} 用字节序列训练，替代 SentencePiece 子词词元序列。
\item \textbf{UL2} 与 T5 类似，采用多种去噪目标进行预训练。
\item \textbf{Flan-T5} 是采用基于提示的 Flan 预训练方法得到的 T5 模型。所用数据集包括 \texttt{taskmaster2}、\texttt{djaym7/wiki\_dialog}、\texttt{deepmind/code\_contests} 等。
\item \textbf{FLan-UL2} 使用 Flan 的提示调优方法和数据集集合微调 UL2。
\item \textbf{UMT5} 是多语言 T5 模型，采用 UniMax 采样，在更新的 mC4 多语言语料库上训练。该语料库含 107 种语言、29 万亿个字符。
\end{enumerate}

% PDF 98--103；原书第 68--73 页。
\subsection{GPT（解码器）}
OpenAI 的生成式预训练 Transformer（Generative Pre-trained Transformer，GPT）系列以较强的语言建模能力改变了自然语言处理领域（\refcite{ch02-brown}{Brown et al., 2020}；\refcite{ch02-radford-2018}{Radford et al., 2018}, \refcite{ch02-radford-2019}{2019}）。值得注意的是，这些模型无需为目标任务进行监督训练，就能完成问答、文本蕴含和文本摘要等自然语言处理任务。本节考察三个关键版本 GPT-1、GPT-2 和 GPT-3，梳理它们的演进过程。GPT-3.5 等后续模型基于采用人类反馈强化学习的 InstructGPT 方法，第 5 章将详细讨论。

\subsubsection{数据集}
\begin{enumerate}
\item GPT-1 最初在 BooksCorpus 数据集上预训练，随后针对多种语言理解任务分别微调。自然语言推断使用 SNLI、MultiNLI、Question NLI、RTE 和 SciTail 等数据集。问答使用 RACE 和 Story Cloze。为衡量语言模型在句子相似度任务上的表现，研究者选用了 MSR Paraphrase Corpus、Quora Question Pairs 和 STS Benchmark 等数据集。分类任务以 Stanford Sentiment Treebank-2 和 CoLA 为基准。
\item 训练 GPT-2 时，作者从 Reddit 上获得较多赞同票的帖子中跟踪外部链接，提取数据，整理成 WebText 数据集。该数据集大小为 40\,GB，含 800 多万份文档，规模大于 GPT-1 使用的 BookCorpus。为保证测试集不受干扰，WebText 排除了维基百科文章。GPT-2 未针对各项任务微调，而是通过零样本推理取得结果。
\item GPT-3 的训练使用五种不同语料，分别赋予特定的采样权重。质量较高的数据集采样更频繁，模型在这些数据上训练多个轮次。所用数据集包括 Common Crawl、WebText2、Books1、Books2 和维基百科。
\end{enumerate}

\subsubsection{架构}
GPT 模型均采用纯解码器架构。表~\ref{tab:ch02-gpt} 列出了不同版本的架构差异。

\sourcefigure{fig02_13.pdf}{GPT-1 的架构及其训练目标。为不同任务微调时，先把结构化输入转换为词元序列，经预训练模型处理，再接线性层与 softmax 层。图右由上到下分别为分类、蕴含、相似度和多项选择的输入结构。Start、Delim 和 Extract 分别标识序列起点、分隔符和用于提取表示的位置。}{fig:ch02-gpt}

\begin{table}[!htbp]
\centering
\caption{GPT-1、GPT-2 和 GPT-3 的比较。}
\label{tab:ch02-gpt}
\small
\begin{tabularx}{\linewidth}{@{}Xrrr@{}}
\toprule
特征 & GPT-1 & GPT-2 & GPT-3\\
\midrule
参数量 & 1.17 亿 & 15 亿 & 1\,750 亿\\
解码器层数 & 12 & 48 & 96\\
上下文词元数 & 512 & 1\,024 & 2\,048\\
隐藏层宽度 & 768 & 1\,600 & 12\,288\\
批量大小 & 64 & 512 & 320 万\\
\bottomrule
\end{tabularx}
\end{table}

\begin{enumerate}
\item GPT-1 使用 12 层的纯解码器 Transformer，结构与原始 Transformer 解码器一致，区别是使用可学习的位置嵌入替代原有的固定位置嵌入，见图~\ref{fig:ch02-gpt}。
\item GPT-2 保留 GPT-1 的架构设计，但规模明显更大。它有 15 亿个参数，比 GPT-1 的 1.17 亿大约多一个数量级。
\item GPT-3 沿用 GPT-1 与 GPT-2 的架构基础，但上下文词元数、层数等与 GPT-2 不同。
\end{enumerate}

\subsubsection{训练}
\begin{enumerate}
\item GPT-1 采用两阶段训练，先无监督预训练，再监督微调。第一阶段在大规模文本语料库上训练容量较大的语言模型。给定无监督词元语料 $U=\{u_1,\ldots,u_n\}$，使用标准的语言建模目标。
\begin{equation}
\mathcal{L}_1(U)=\sum_i\log P(u_i\mid u_{i-k},\ldots,u_{i-1};\Theta),
\tag{2.34}
\end{equation}
其中 $k$ 是上下文窗口大小，条件概率 $P$ 由参数为 $\Theta$ 的神经网络建模。这些参数通过随机梯度下降（stochastic gradient descent）训练。

无监督预训练后，将模型参数适配到一个监督目标任务。给定有标签数据集 $C$，每个样例由输入词元序列 $x_1,\ldots,x_m$ 和标签 $y$ 组成。输入经过预训练模型，得到最后一个 Transformer 块的激活值 $hm_l$。
\begin{equation}
P(y\mid x_1,\ldots,x_m)=\operatorname{softmax}(hm_lW_y).
\tag{2.35}
\end{equation}
相应的目标为
\begin{equation}
\mathcal{L}_2(C)=\sum_{(x,y)}\log P(y\mid x_1,\ldots,x_m).
\tag{2.36}
\end{equation}
此外，在微调中把语言建模作为辅助目标，可以改善泛化和收敛。合成目标为
\begin{equation}
\mathcal{L}_3=\mathcal{L}_2(C)+\lambda\times\mathcal{L}_1(U),
\tag{2.37}
\end{equation}
其中 $\mathcal{L}_1(U)$ 是无监督预训练目标，$\mathcal{L}_2(C)$ 是监督微调目标。合成目标 $\mathcal{L}_3$ 同时利用这两个阶段的目标。针对任务转换输入，使预训练模型能够处理多种任务的结构化输入，无须明显改动架构。

\item GPT-2 的主要训练方法基于语言建模。可把它看作从一组样例 $(x_1,x_2,\ldots,x_n)$ 中进行无监督分布估计，其中各 $x_i$ 是一个符号序列。模型以输入和具体任务为条件，以便泛化到多种任务。形式上，模型旨在估计
\begin{equation}
p(\text{output}\mid\text{input},\text{task}).
\tag{2.38}
\end{equation}
例如，翻译任务可在模型中表示为 \texttt{("translate to French", "English text", "French text")}。以往语言模型通常在单一领域上训练，而这里的方法强调使用涵盖多个领域与语境的多样化数据集，以捕捉数据集一节所述的多种自然语言模式。

\item GPT-3 在上文所述的大规模语料库上，以自回归方式预测下一个词来训练。它用上下文学习改变了传统上利用专用训练数据针对特定任务微调的做法。GPT-3 可以只凭提示中的上下文或示例，动态适应未明确训练过的新任务。不同的学习方式进一步体现了 GPT-3 上下文学习的灵活性。在少样本学习中，提示包含多个示例来引导模型。例如，要将英语译成法语，可以先给出几对英法句子，再给出需要翻译的英语句子。单样本学习（one-shot learning）只给一个引导示例。零样本学习没有明确示例，GPT-3 根据描述性提示执行任务，展示出单凭预训练就能理解并完成任务的能力。接下来几章将深入讨论这个主题。
\end{enumerate}

\subsubsection{要点}
\begin{keypoint}
\begin{enumerate}
\item GPT-1 证明了以语言建模作预训练的效果，模型由此具备较强的泛化能力。其架构支持迁移学习，只需少量微调即可处理多种自然语言处理任务。这个模型显示了生成式预训练的潜力，为后续模型利用更大的数据集和更多参数来发挥这些能力开辟了道路。
\item GPT-2 证明了增加数据量和参数量可以增强语言模型的能力，在零样本情形中常常超过当时的先进结果。研究发现，表现随模型容量呈对数线性增长。模型的困惑度（perplexity）随参数量增加而持续下降，没有出现饱和迹象。GPT-2 对 WebText 数据欠拟合，这说明继续训练可能改善结果。该发现表明，进一步增大模型或能改善自然语言理解。
\item 在零样本设置中，GPT-2 在 8 个语言建模数据集中的 7 个上超过当时的基准结果。在 Children’s Book 数据集上，普通名词和命名实体识别的准确率比当时的先进结果高约 7\%。在 LAMBADA 上，GPT-2 把困惑度从 99.8 显著降到 8.6，表明准确率明显提高。在阅读理解任务中，GPT-2 超过了四个基准模型中的三个。但在文本摘要方面，它的表现与传统上专门训练的摘要模型相当，甚至更低。
\item 研究者在多个语言建模和自然语言处理数据集上评估了 GPT-3。它在 LAMBADA 和 Penn Tree Bank 等数据集上表现出色，经常达到或超过先进模型，少样本或零样本设置尤其如此。它并非总能超过最好的基准结果，但零样本表现持续改善。GPT-3 擅长多种自然语言处理任务，包括闭卷问答（closed-book question answering）与翻译，表现往往可以与微调模型相比或超过它们。它在少样本场景中通常比单样本或零样本场景表现更好。研究还用算术、打乱字母复原单词和新词用法等不常规任务测试 GPT-3。它在这些任务上的表现随参数量扩大而改善，在少样本设置下尤为明显。
\item GPT-3 能生成高质量文本，但较长句子有时不够连贯，文本也容易重复。它难以完成自然语言推断、填空及部分阅读理解任务，原因可能是其单向结构。研究指出，双向模型将来或能解决这个问题。GPT-3 的目标对所有词元一视同仁，缺少面向特定任务的预测。研究讨论了扩充目标、强化学习或加入其他模态等做法，也指出 GPT-3 的架构规模大，推理复杂且昂贵，输出难以解释。此外，它能生成类似人类写作的文本，这也带来了被滥用于网络钓鱼或散布虚假信息的风险。
\end{enumerate}
\end{keypoint}

\subsubsection{变体}
\begin{enumerate}
\item \textbf{Gopher。}这个模型有 2\,800 亿个参数，使用 MassiveText 数据集中的 3\,000 亿个词元训练，上下文窗口为 2\,048 个词元。数据集包含网页、图书、新闻和代码。在 100 项任务中，Gopher 有 81\% 的任务超过 GPT-3（1\,750 亿个参数）等当时的先进模型（\refcite{ch02-rae-2021}{Rae et al., 2021}）。
\item \textbf{Chinchilla。}这个按算力最优原则训练的模型有 700 亿个参数，训练使用了 1.4 万亿个词元。它的表现超过规模更大的 Gopher，而且模型更小，显著降低了推理成本（\refcite{ch02-hoffmann}{Hoffmann et al., 2022}）。
\item \textbf{Llama。}Meta 的 GPT 类 Llama 模型在本书写作时已到第 3 版，是一种有 80 亿和 700 亿参数两个版本的开源大语言模型。它提供预训练和指令调优模型，使用监督微调和基于人类反馈的强化学习来优化对话能力（\refcite{ch02-touvron}{Touvron et al., 2023}）。
\item \textbf{Claude。}Anthropic 的 Claude 3 系列包括 Opus、Sonnet 和 Haiku，原书所列规模分别为 200 亿、700 亿和 2 万亿个参数，各自针对不同的性能需求（\refcite{ch02-anthropic}{Anthropic, 2023}）。它们支持多语言和视觉处理，且更易受指令引导。Opus 在复杂任务上提供最高水平的表现，Sonnet 兼顾表现和成本，Haiku 则最快且最便宜。当提示长度低于 32\,000 个词元，输入与输出词元之比为 1:5 时，Haiku 每秒处理 21\,000 个词元。
\item \textbf{Command R。}Cohere 开发的 Command R 有 350 亿个参数，是面向检索增强生成、调用外部 API 和工具等长上下文任务优化的生成模型。它面向规模化应用设计，在检索增强生成和工具使用中有较高的准确率，延迟低、吞吐量高，上下文长度为 128\,000 个词元（\refcite{ch02-cohere}{Cohere, 2024}）。Command R 在 10 种主要语言中也表现出较强的能力。
\item \textbf{Gemma。}Google 开发的 Gemma 系列包括 20 亿与 70 亿参数两个版本，使用 6 万亿个词元训练（\refcite{ch02-team}{Team et al., 2024}）。这些模型在语言理解、推理和安全方面表现突出，在 18 项任务中的 11 项超过规模相当的模型。
\end{enumerate}

% PDF 103--106；原书第 73--76 页。
\subsection{Mixtral 8x7B（混合专家）}
自 \refcite{ch02-jacobs}{Jacobs 等人（1991）}提出混合专家模型以来，这类模型已有很大发展。最初，研究者用混合专家把复杂问题分成较容易处理的子问题，合并多个专家网络的输出。各个网络分别擅长整体任务的不同方面，门控网络（gating network）根据输入动态决定各专家的贡献权重。

Shazeer 等人于 2017 年引入 Top-$k$ 路由（top-$k$ routing），是一次重要进展。此法只计算排名最前的 $k$ 个专家的输出，使模型可以拥有数十亿个参数，同时让计算成本保持在可承受的范围内，并在语言建模等任务上显著改善表现。

Switch Transformer 又把 Top-$k$ 路由推进了一步，采用“硬路由”（hard routing），令 $k=1$，每个输入词元只选择一个最相关的专家（\refcite{ch02-lepikhin}{Lepikhin et al., 2020}）。这个模型用 128 个采取硬路由的专家取代 T5 Transformer 中传统的前馈网络层，并结合多种优化技术，提高训练效率以及在 GLUE 基准等任务上的表现。

Mixtral 8x7B 是一个高质量的稀疏混合专家模型（Sparse Mixture of Experts Model，SMoE），以 Apache 2.0 许可证公开发布（\refcite{ch02-jiang}{Jiang et al., 2024}）。它在多数基准上的表现超过 Llama-2 70B，推理速度达到后者的 6 倍，并在多数标准基准上达到或超过 GPT-3.5。

\subsubsection{数据集}
原书未列出预训练细节。据报道，该模型使用从开放网页语料库收集的多语言数据集训练。它可以处理多种语言，包括英语、法语、意大利语、德语和西班牙语。

\subsubsection{架构}
\begin{figure}[!htbp]
\centering
% 重绘自原书第 74 页的图 2.14。
\begin{tikzpicture}[
  font=\small\sffamily,
  stage/.style={draw=orange!75!black,fill=yellow!32,rounded corners=2mm,
    minimum width=18mm,minimum height=9mm,inner sep=1pt},
  wire/.style={draw=black!60,line width=.55pt},
  flow/.style={wire,-{Stealth[length=1.5mm]}}]
\node (inputs) at (.1,0) {输入};
\node[stage] (router) at (1.95,0) {路由器};
\foreach \x/\y in {5.05/.62,5.31/.43,5.57/.24,5.83/.05} {
  \node[stage,draw=orange!52!gray,densely dotted] at (\x,\y) {};
}
\node[stage,fill=yellow!29] (rear) at (5.17,.50) {};
\node[stage] (front) at (6.09,-.16) {专家};
\node[circle,draw=black!65,minimum size=5.4mm,inner sep=0pt] (sum) at (8.20,.03) {$+$};
\node (outputs) at (9.65,.03) {输出};

\draw[flow] (inputs.east) -- (router.west);
\draw[wire] (router.east) -- (3.65,0);
\draw[flow] (3.65,0) |- (rear.west);
\draw[flow] (3.65,0) -- (3.93,0) |- (front.west);
\draw[flow] (rear.east) -- (7.33,.50) |- (sum.west);
\draw[flow] (front.east) -- (7.05,-.16) |- (8.20,-.55) -- (sum.south);
\draw[wire] (router.north) -- (1.95,1.76) -- (7.72,1.76);
\node[anchor=north west] at (2.02,1.76) {门控权重};
\draw[flow] (7.46,1.76) -- (7.46,.03);
\draw[flow] (7.72,1.76) -- (7.72,-.55);
\draw[flow] (sum.east) -- (outputs.west);
\end{tikzpicture}

\caption{Mixtral 的混合专家层。每个输入向量被路由到 8 个专家中的 2 个。该层输出是被选中专家的输出的加权和。专家采用传统 Transformer 架构中的标准前馈模块。}
\label{fig:ch02-mixtral}
\end{figure}

如图~\ref{fig:ch02-mixtral} 所示，Mixtral 将传统 Transformer 的标准前馈网络子模块换成混合专家层，每个词元独立处理。对于输入向量 $\symbf{x}$，混合专家层的输出是多个专家网络输出的加权和。权重由门控网络的输出决定。设 $n$ 个专家网络为 $\{E_0,E_1,\ldots,E_{n-1}\}$，则混合专家层的输出为
\begin{equation}
y=\sum_{i=0}^{n-1}G(\symbf{x})_i\cdot E_i(\symbf{x}),
\tag{2.39}
\end{equation}
其中 $G(\symbf{x})_i$ 是门控网络对第 $i$ 个专家的输出，$E_i(\symbf{x})$ 是第 $i$ 个专家网络的输出。若门控向量是稀疏的，则可以跳过权重为零的专家。

$G(\symbf{x})$ 有多种实现（\refcite{ch02-clark}{Clark et al., 2022}；\refcite{ch02-hazimeh}{Hazimeh et al., 2021}）。一种简单而高效的方法是，取线性层输出中最大的 $k$ 个 logit，再对它们计算 softmax（\refcite{ch02-shazeer-2017}{Shazeer et al., 2017}）。门控函数定义为
\begin{equation}
G(\symbf{x}):=\operatorname{Softmax}(\operatorname{TopK}(\symbf{x}\cdot\symbf{W}_g)),
\tag{2.40}
\end{equation}
其中，若 $\ell_i$ 是 logit 向量 $\ell\in\mathbb{R}^{n}$ 中最大的 $k$ 个分量之一，则 $\operatorname{TopK}(\ell)_i=\ell_i$，否则 $\operatorname{TopK}(\ell)_i=-\infty$。$k$ 表示每个词元使用的专家数，是控制每词元计算量的超参数。

Mixtral 使用 SwiGLU 架构实现专家函数 $E_i(\symbf{x})$，并取 $k=2$。因此，每个词元会被送到两个权重不同的 SwiGLU 子模块。给定输入词元 $\symbf{x}$，输出 $y$ 的计算方式为
\begin{equation}
y=\sum_{i=0}^{n-1}\operatorname{Softmax}(\operatorname{Top2}(\symbf{x}\cdot\symbf{W}_g))_i
\cdot\operatorname{SwiGLU}_i(\symbf{x}).
\tag{2.41}
\end{equation}

\subsubsection{训练}
研究者没有提供 Mixtral 8x7B 的数据预处理、训练方法或训练硬件信息。

\subsubsection{要点}
\begin{keypoint}
\begin{enumerate}
\item 研究发现，在预训练中大幅提高多语言数据的比例，使 Mixtral 在多语言基准上表现出色，同时保持良好的英语表现。与 Llama-2 70B 相比，Mixtral 在法语、德语、西班牙语和意大利语上表现明显更好。结果表明，在 ARC Challenge、HellaSwag 和 MMLU 等基准上，Mixtral 对这些语言的表现优于 Llama-2 70B。
\item 使用专用的高性能计算内核，可以在单个 GPU 上高效执行混合专家层。MegaBlocks（\refcite{ch02-gale}{Gale et al., 2023}）提供的内核就是一例，它将混合专家层中的前馈网络运算作为大规模稀疏矩阵乘法来处理。也可以把混合专家层有效分布到多个 GPU，既使用标准并行方式，又使用称为专家并行（Expert Parallelism，EP）的专门分片方法（\refcite{ch02-shazeer-2017}{Shazeer et al., 2017}）。执行时，分给某位专家的词元由对应的 GPU 处理，输出随后被送回该词元在序列中的原位置。
\item 这种设置需要大量显存（Video Random Access Memory，VRAM）。即使在某一时刻只有一两个专家参与运算，所有专家也必须载入显存。
\end{enumerate}
\end{keypoint}

\subsubsection{变体}
\begin{enumerate}
\item \textbf{Mixtral 8x22B} 是规模较大的稀疏混合专家变体，共有至多 1\,410 亿个参数，推理时使用约 390 亿个。它提高了推理吞吐量，但也需要更多显存。这个模型可处理至多 64\,000 个词元。
\item \textbf{Mixtral 8x7B--Instruct} 面向聊天，使用成对反馈数据进行监督微调和直接偏好优化，研究者没有公布具体数据集。LMSys 发布的人工评估结果表明，这一版本的 Elo 评分高于 GPT-3.5 Turbo 和 Claude 2.1。
\item \textbf{MegaBlocks} 通过块稀疏运算（block-sparse operation）和定制的 GPU 内核，提高混合专家模型在 GPU 上的训练效率，改善词元利用率与硬件映射。它的训练速度比 Tutel 程序库最多快 40\%，可达 Megatron-LM 框架训练稠密神经网络的 2.4 倍，同时保持模型质量。
\end{enumerate}

% PDF 106--108；原书第 76--78 页。
\section[教程：认识大语言模型与预训练]{教程——认识大语言模型与预训练}
\subsection{概述}
本教程考察大语言模型架构如何工作，重点比较掩码模型与因果模型。先考察现有的预训练模型，了解它们怎样产生输出。看清大语言模型发挥作用的方式后，再运行一个简化的训练循环，初步了解训练过程。

\noindent\textbf{目标}
\begin{itemize}
\item 检查大语言模型的输入和输出，包括词元切分器。
\item 逐步查看代码，演示掩码大语言模型和因果大语言模型如何预测词元。
\item 以小规模例子演示怎样从头训练大语言模型。
\item 检查训练循环是否按预期工作。
\end{itemize}
这里给出的是教程的简版。完整版可在 \url{https://github.com/springer-llms-deep-dive/llms-deep-dive-tutorials} 查阅。

\subsection{实验设计}
本教程最终要展示预训练过程怎样运作。我们先分析大语言模型架构的组成部分。首先考察前向传播（forward pass），借此介绍各部件以及它们如何协作，实现语言建模目标。BERT 与 GPT-2 都使用这段代码，以对比掩码模型（纯编码器）与自回归模型（纯解码器）的异同。

分析完词元预测的步骤后，就容易理解大语言模型的训练循环。它采用典型的反向传播（backpropagation），使梯度经过模型的各层。这里假定读者了解深度学习基础，不再讨论具体超参数的影响或训练循环的其他细节。需要简要复习的读者可以参阅附录。

练习结束后，代码会得到一个记住一小部分维基百科数据的示例模型。随书提供的笔记本只包含 GPT-2 的训练循环。读者有需要时，可以把这个实验扩展到其他大语言模型，包括掩码模型。

\subsection{结果与分析}
在大语言模型预训练实验中，训练损失很快下降，验证损失却保持在高位，如图~\ref{fig:ch02-loss} 所示。这是模型过拟合训练数据时会出现的现象。要使模型习得足够多的信息，在验证数据上表现出良好的泛化能力，需要更多文档和训练步骤。英语中可能出现的词元序列数量庞大，出现这种结果并不奇怪。

\sourcefigure{fig02_15.pdf}{GPT-2 学习一小组维基百科文档内容时的损失曲线。横轴为训练步数，纵轴为损失，蓝线为训练集，橙线为验证集。}{fig:ch02-loss}

虽然训练还不足以让模型在验证数据上表现良好，可以看出它已从训练数据中学到了很多内容。可以用一个训练样例来验证。
\begin{lstlisting}[language=Python,caption={访问数据集中的文本样例}]
print(raw_datasets["train"][0]["text"])
# Output:
# William Edward Whitehouse (20 May 1859 – 12 January 1935) was an English cellist.
# Career
# He studied for one year with Alfredo Piatti, for whom he deputised (taking his place in concerts when called upon), and was his favourite pupil. He went on to teach at the Royal Academy of Music, Royal College of Music and King's College, Cambridge...
\end{lstlisting}

给定开头若干词元，可以检查模型是否能补全这段文本。
\begin{lstlisting}[language=Python,caption={用模型生成文本}]
text = "William Edward Whitehouse (20 May 1859 – 12 January 1935) was an English cellist.\n\nCareer\nHe studied for one year with"

model_inputs = tokenizer(text, return_tensors='pt')
output_generate = model.generate(**model_inputs, max_new_tokens=5)
sequence = tokenizer.decode(output_generate[0])
print(sequence)
# Output:
# William Edward Whitehouse (20 May 1859 – 12 January 1935) was an English cellist.
#
# Career
# He studied for one year with Alfredo Piatti,
\end{lstlisting}

在这个例子中，模型正确识别了 Alfredo Piatti，说明它反复接触特定的维基百科条目后记住了这条信息。这让我们相信，词元切分器与模型能从维基百科中学习语言模式。当然，这并不保证同样的训练方法可以直接用于完整规模的数据集。学习率等具体参数可能需要调整。

\subsection{结论}
我们展示了掩码语言模型与因果语言模型如何预测词元，并演示了这些模型如何反复尝试预测词元，随后更新权重以降低损失，由此把信息学进模型。

% PDF 109--112；原书第 79--82 页。
\begin{chapterreferences}
\urldef\hochreiterurl\url{http://dblp.uni-trier.de/db/journals/ijufks/ijufks6.html#Hochreiter98}
\refentry{ch02-anthropic}{Anthropic. The Claude 3 model family, 2023. URL \url{https://www.anthropic.com/news/claude-3-family}.}
\refentry{ch02-ba}{Jimmy Lei Ba, Jamie Ryan Kiros, and Geoffrey E Hinton. Layer normalization. \textit{arXiv preprint arXiv:1607.06450}, 2016.}
\refentry{ch02-bahdanau}{Dzmitry Bahdanau, Kyunghyun Cho, and Yoshua Bengio. Neural machine translation by jointly learning to align and translate. CoRR, abs/1409.0473, 2014.}
\refentry{ch02-biderman}{Stella Biderman et al. Pythia: A suite for analyzing large language models across training and scaling, 2023.}
\refentry{ch02-brown}{Tom Brown, Benjamin Mann, Nick Ryder, Melanie Subbiah, Jared D Kaplan, Prafulla Dhariwal, Arvind Neelakantan, Pranav Shyam, Girish Sastry, Amanda Askell, et al. Language models are few-shot learners. \textit{Advances in Neural Information Processing Systems}, 33:1877–1901, 2020.}
\refentry{ch02-clark}{Aidan Clark, Diego de Las Casas, Aurelia Guy, Arthur Mensch, Michela Paganini, Jordan Hoffmann, Bogdan Damoc, Blake Hechtman, Trevor Cai, Sebastian Borgeaud, et al. Unified scaling laws for routed language models. In \textit{International Conference on Machine Learning}, pages 4057–4086. PMLR, 2022.}
\refentry{ch02-cohere}{Cohere. Command R: A generative model for long context tasks, 2024. URL \url{https://cohere.com/command}.}
\refentry{ch02-devlin}{Jacob Devlin, Ming-Wei Chang, Kenton Lee, and Kristina Toutanova. BERT: Pre-training of deep bidirectional transformers for language understanding, 2019.}
\refentry{ch02-ding}{Ming Ding, Zhuoyi Yang, Wenyi Hong, Wendi Zheng, Chang Zhou, Da Yin, Junyang Lin, Xu Zou, Zhou Shao, Hongxia Yang, et al. CogView: Mastering text-to-image generation via transformers. \textit{Advances in Neural Information Processing Systems}, 34:19822–19835, 2021.}
\refentry{ch02-fu}{Hao Fu, Yao; Peng and Tushar Khot. How does GPT obtain its ability? Tracing emergent abilities of language models to their sources. Yao Fu’s Notion, Dec 2022. URL \url{https://yaofu.notion.site/How-does-GPT-Obtain-its-Ability-Tracing-Emergent-Abilities-of-Language-Models-to-their-Sources-b9a57ac0fcf74f30a1ab9e3e36fa1dc1}.}
\refentry{ch02-gale}{Trevor Gale, Deepak Narayanan, Cliff Young, and Matei Zaharia. MegaBlocks: Efficient sparse training with mixture-of-experts. \textit{Proceedings of Machine Learning and Systems}, 5, 2023.}
\refentry{ch02-gao}{Leo Gao et al. The Pile: An 800GB dataset of diverse text for language modeling, 2020.}
\refentry{ch02-harlap}{Aaron Harlap, Deepak Narayanan, Amar Phanishayee, Vivek Seshadri, Nikhil Devanur, Greg Ganger, and Phil Gibbons. PipeDream: Fast and efficient pipeline parallel DNN training. \textit{arXiv preprint arXiv:1806.03377}, 2018.}
\refentry{ch02-hazimeh}{Hussein Hazimeh, Zhe Zhao, Aakanksha Chowdhery, Maheswaran Sathiamoorthy, Yihua Chen, Rahul Mazumder, Lichan Hong, and Ed Chi. DSelect-k: Differentiable selection in the mixture of experts with applications to multi-task learning. \textit{Advances in Neural Information Processing Systems}, 34:29335–29347, 2021.}
\refentry{ch02-he}{Ji He, Jianshu Chen, Xiaodong He, Jianfeng Gao, Lihong Li, Li Deng, and Mari Ostendorf. Deep reinforcement learning with a natural language action space, 2016.}
\refentry{ch02-hernandez}{Danny Hernandez et al. Scaling laws and interpretability of learning from repeated data, 2022.}
\refentry{ch02-hochreiter}{Sepp Hochreiter. The vanishing gradient problem during learning recurrent neural nets and problem solutions. \textit{International Journal of Uncertainty, Fuzziness and Knowledge-Based Systems}, 6(2):107–116, 1998. URL \hochreiterurl.}
\refentry{ch02-hoffmann}{Jordan Hoffmann et al. Training compute-optimal large language models, 2022.}
\refentry{ch02-huang}{Yanping Huang, Youlong Cheng, Ankur Bapna, Orhan Firat, Dehao Chen, Mia Chen, HyoukJoong Lee, Jiquan Ngiam, Quoc V Le, Yonghui Wu, et al. GPipe: Efficient training of giant neural networks using pipeline parallelism. \textit{Advances in Neural Information Processing Systems}, 32, 2019.}
\refentry{ch02-jacobs}{Robert A Jacobs, Michael I Jordan, Steven J Nowlan, and Geoffrey E Hinton. Adaptive mixtures of local experts. \textit{Neural Computation}, 3(1):79–87, 1991.}
\refentry{ch02-jiang}{Albert Q Jiang, Alexandre Sablayrolles, Antoine Roux, Arthur Mensch, Blanche Savary, Chris Bamford, Devendra Singh Chaplot, Diego de las Casas, Emma Bou Hanna, Florian Bressand, et al. Mixtral of experts. \textit{arXiv preprint arXiv:2401.04088}, 2024.}
\refentry{ch02-kalyan}{Katikapalli Subramanyam Kalyan, Ajit Rajasekharan, and Sivanesan Sangeetha. AMMUS: A survey of transformer-based pretrained models in natural language processing. \textit{arXiv preprint arXiv:2108.05542}, 2021.}
\refentry{ch02-kamath}{Uday Kamath, Kenneth L Graham, and Wael Emara. \textit{Transformers for Machine Learning: A Deep Dive}. CRC Press, 2022.}
\refentry{ch02-laurencon}{Hugo Laurençon et al. The BigScience ROOTS corpus: A 1.6TB composite multilingual dataset, 2023.}
\refentry{ch02-lepikhin}{Dmitry Lepikhin, HyoukJoong Lee, Yuanzhong Xu, Dehao Chen, Orhan Firat, Yanping Huang, Maxim Krikun, Noam Shazeer, and Zhifeng Chen. GShard: Scaling giant models with conditional computation and automatic sharding. \textit{arXiv preprint arXiv:2006.16668}, 2020.}
\refentry{ch02-liello}{Luca Di Liello, Matteo Gabburo, and Alessandro Moschitti. Efficient pre-training objectives for transformers, 2021.}
\refentry{ch02-lin}{Tianyang Lin, Yuxin Wang, Xiangyang Liu, and Xipeng Qiu. A survey of transformers. \textit{AI Open}, 2022.}
\refentry{ch02-luong}{Minh-Thang Luong, Hieu Pham, and Christopher D. Manning. Effective approaches to attention-based neural machine translation, 2015.}
\refentry{ch02-nair}{Vinod Nair and Geoffrey E Hinton. Rectified linear units improve restricted Boltzmann machines. In \textit{Proceedings of the 27th International Conference on Machine Learning (ICML-10)}, pages 807–814, 2010.}
\refentry{ch02-press}{Ofir Press, Noah A Smith, and Mike Lewis. Train short, test long: Attention with linear biases enables input length extrapolation. \textit{arXiv preprint arXiv:2108.12409}, 2021.}
\refentry{ch02-radford-2018}{Alec Radford, Karthik Narasimhan, Tim Salimans, Ilya Sutskever, et al. Improving language understanding by generative pre-training. 2018.}
\refentry{ch02-radford-2019}{Alec Radford, Jeffrey Wu, Rewon Child, David Luan, Dario Amodei, Ilya Sutskever, et al. Language models are unsupervised multitask learners. \textit{OpenAI blog}, 1(8):9, 2019.}
\refentry{ch02-rae-2021}{Jack W Rae et al. Scaling language models: Methods, analysis \& insights from training Gopher. \textit{arXiv preprint arXiv:2112.11446}, 2021.}
\refentry{ch02-rae-2022}{Jack W. Rae et al. Scaling language models: Methods, analysis insights from training Gopher, 2022.}
\refentry{ch02-raffel}{Colin Raffel, Noam Shazeer, Adam Roberts, Katherine Lee, Sharan Narang, Michael Matena, Yanqi Zhou, Wei Li, and Peter J. Liu. Exploring the limits of transfer learning with a unified text-to-text Transformer, 2020.}
\refentry{ch02-rajbhandari}{Samyam Rajbhandari, Jeff Rasley, Olatunji Ruwase, and Yuxiong He. ZeRO: Memory optimizations toward training trillion parameter models. In \textit{SC20: International Conference for High Performance Computing, Networking, Storage and Analysis}, pages 1–16. IEEE, 2020.}
\refentry{ch02-ramachandran}{Prajit Ramachandran, Barret Zoph, and Quoc V Le. Searching for activation functions. \textit{arXiv preprint arXiv:1710.05941}, 2017.}
\refentry{ch02-shazeer-2020}{Noam Shazeer. GLU variants improve Transformer. \textit{arXiv preprint arXiv:\allowbreak{}2002.05202}, 2020.}
\refentry{ch02-shazeer-2017}{Noam Shazeer, Azalia Mirhoseini, Krzysztof Maziarz, Andy Davis, Quoc Le, Geoffrey Hinton, and Jeff Dean. Outrageously large neural networks: The sparsely-gated mixture-of-experts layer. \textit{arXiv preprint arXiv:1701.06538}, 2017.}
\refentry{ch02-shoeybi}{Mohammad Shoeybi, Mostofa Patwary, Raul Puri, Patrick LeGresley, Jared Casper, and Bryan Catanzaro. Megatron-LM: Training multi-billion parameter language models using model parallelism. \textit{arXiv preprint arXiv:1909.08053}, 2019.}
\refentry{ch02-su}{Jianlin Su, Yu Lu, Shengfeng Pan, Ahmed Murtadha, Bo Wen, and Yunfeng Liu. RoFormer: Enhanced Transformer with rotary position embedding. \textit{arXiv preprint arXiv:2104.09864}, 2021.}
\refentry{ch02-sutskever}{Ilya Sutskever, Oriol Vinyals, and Quoc V. Le. Sequence to sequence learning with neural networks, 2014.}
\refentry{ch02-team}{Gemma Team, Mesnard, et al. Gemma: Open models based on Gemini research and technology. \textit{arXiv preprint arXiv:2403.08295}, 2024.}
\refentry{ch02-touvron}{Hugo Touvron et al. Llama 2: Open foundation and fine-tuned chat models, 2023.}
\refentry{ch02-vaswani}{Ashish Vaswani, Noam Shazeer, Niki Parmar, Jakob Uszkoreit, Llion Jones, Aidan N Gomez, Łukasz Kaiser, and Illia Polosukhin. Attention is all you need. In I. Guyon, U. Von Luxburg, S. Bengio, H. Wallach, R. Fergus, S. Vishwanathan, and R. Garnett, editors, \textit{Advances in Neural Information Processing Systems}, volume 30. Curran Associates, Inc., 2017. URL \url{https://proceedings.neurips.cc/paper_files/paper/2017/file/3f5ee243547dee91fbd053c1c4a845aa-Paper.pdf}.}
\refentry{ch02-wang-2019}{Alex Wang, Amanpreet Singh, Julian Michael, Felix Hill, Omer Levy, and Samuel R. Bowman. GLUE: A multi-task benchmark and analysis platform for natural language understanding, 2019.}
\refentry{ch02-wang-2022}{Hongyu Wang, Shuming Ma, Li Dong, Shaohan Huang, Dongdong Zhang, and Furu Wei. DeepNet: Scaling transformers to 1,000 layers. \textit{arXiv preprint arXiv:2203.00555}, 2022.}
\refentry{ch02-workshop}{BigScience Workshop et al. BLOOM: A 176B-parameter open-access multilingual language model, 2023.}
\refentry{ch02-zhang}{Biao Zhang and Rico Sennrich. Root mean square layer normalization. \textit{Advances in Neural Information Processing Systems}, 32, 2019.}
\refentry{ch02-zhao}{Wayne Xin Zhao et al. A survey of large language models, 2023.}
\end{chapterreferences}

